\section{2023年高考题}
\subsection{2023年全国甲卷（理）}\label{gk:2023年全国甲卷（理）}

\begin{description}
    \item[一、选择题] 本题共12小题，每小题5分，共60分．在每小题给出的四个选项中，只有一项是符合题目要求的．
    \begin{enumerate}[leftmargin=5pt]
        \item 设集合 $A=\{x\mid x=3k+1,k\in\mathbb{Z}\}$，$B=\{x\mid x=3k+2,k\in\mathbb{Z}\}$，$U$ 为整数集，则 $\complement_U(A\cup B)=$
       
        {\centering\begin{tabular*}{\linewidth}{@{}@{\extracolsep{\fill}}llll@{}}
          \(\text{A. }\{x\mid x=3k,k\in\mathbb{Z}\}\)  & \(\text{B. }\{x\mid x=3k-1,k\in\mathbb{Z}\}\)  & \(\text{C. }\{x\mid x=3k-2,k\in\mathbb{Z}\}\)  & \(\text{D. }\varnothing\)
        \end{tabular*}\par}

        \item 若复数 $(a+\i)(1-a\i)=2$，则 $a=$
       
        {\centering\begin{tabular*}{\linewidth}{@{}@{\extracolsep{\fill}}llll@{}}
          \(\text{A. }-1\)  & \(\text{B. }0\)  & \(\text{C. }1\)  & \(\text{D. }2\)
        \end{tabular*}\par}

        \item 执行下面的程序框图，输出的 $B=$
        
        \textsf{[插入图片：程序框图]}

        {\centering\begin{tabular*}{\linewidth}{@{}@{\extracolsep{\fill}}llll@{}}
          \(\text{A. }21\)  & \(\text{B. }34\)  & \(\text{C. }55\)  & \(\text{D. }89\)
        \end{tabular*}\par}

        \item 已知向量 $\bm a,\bm b,\bm c$ 满足 $|\bm a|=|\bm b|=1$，$|\bm c|=\sqrt{2}$，且 $\bm a+\bm b+\bm c=\bm 0$，则 $\cos\langle\bm a-\bm c,\bm b-\bm c\rangle=$
       
        {\centering\begin{tabular*}{\linewidth}{@{}@{\extracolsep{\fill}}llll@{}}
          \(\text{A. }-\dfrac{4}{5}\)  & \(\text{B. }-\dfrac{2}{5}\)  & \(\text{C. }\dfrac{2}{5}\)  & \(\text{D. }\dfrac{4}{5}\)
        \end{tabular*}\par}

        \item 设等比数列 $\{a_n\}$ 的各项均为正数，前 $n$ 项和为 $S_n$，若 $a_1=1$，$S_5=5S_3-4$，则 $S_4=$
       
        {\centering\begin{tabular*}{\linewidth}{@{}@{\extracolsep{\fill}}llll@{}}
          \(\text{A. }\dfrac{15}{8}\)  & \(\text{B. }\dfrac{65}{8}\)  & \(\text{C. }15\)  & \(\text{D. }40\)
        \end{tabular*}\par}

        \item 某地的中学生中有 $60\%$ 的同学爱好滑冰，$50\%$ 的同学爱好滑雪，$70\%$ 的同学爱好滑冰或爱好滑雪．在该地的中学生中随机调查一位同学，若该同学爱好滑雪，则该同学也爱好滑冰的概率为
       
        {\centering\begin{tabular*}{\linewidth}{@{}@{\extracolsep{\fill}}llll@{}}
          \(\text{A. }0.8\)  & \(\text{B. }0.6\)  & \(\text{C. }0.5\)  & \(\text{D. }0.4\)
        \end{tabular*}\par}

        \item 设甲：$\sin^2\alpha+\sin^2\beta=1$，乙：$\sin\alpha+\cos\beta=0$，则
       
        {\centering\begin{tabular*}{\linewidth}{@{}@{\extracolsep{\fill}}llll@{}}
          \(\text{A. }\)甲是乙的充分条件但不是必要条件  & \(\text{B. }\)甲是乙的必要条件但不是充分条件 \\
          \(\text{C. }\)甲是乙的充要条件  & \(\text{D. }\)甲既不是乙的充分条件也不是乙的必要条件
        \end{tabular*}\par}

        \item 已知双曲线 $\dfrac{x^2}{a^2}-\dfrac{y^2}{b^2}=1\,(a>0,b>0)$ 的离心率为 $\sqrt{5}$，其中一条渐近线与圆 $(x-2)^2+(y-3)^2=1$ 交于 $A,B$ 两点，则 $|AB|=$
       
        {\centering\begin{tabular*}{\linewidth}{@{}@{\extracolsep{\fill}}llll@{}}
          \(\text{A. }\dfrac{1}{5}\)  & \(\text{B. }\dfrac{\sqrt{5}}{5}\)  & \(\text{C. }\dfrac{2\sqrt{5}}{5}\)  & \(\text{D. }\dfrac{4\sqrt{5}}{5}\)
        \end{tabular*}\par}

        \item 有五名志愿者参加社区服务，共服务星期六、星期天两天，每天从中任选三人参加服务，则恰有一人连续参加两天服务的概率为
       
        {\centering\begin{tabular*}{\linewidth}{@{}@{\extracolsep{\fill}}llll@{}}
          \(\text{A. }\dfrac{1}{5}\)  & \(\text{B. }\dfrac{2}{5}\)  & \(\text{C. }\dfrac{3}{5}\)  & \(\text{D. }\dfrac{4}{5}\)
        \end{tabular*}\par}

        \item 已知 $f(x)$ 为函数 $y=\cos\bigl(2x+\frac{\pi}{6}\bigr)$ 向左平移 $\frac{\pi}{6}$ 个单位所得函数，$y=f(x)$ 与 $y=\frac{1}{2}x-\frac{1}{2}$ 的交点个数为
       
        {\centering\begin{tabular*}{\linewidth}{@{}@{\extracolsep{\fill}}llll@{}}
          \(\text{A. }1\)  & \(\text{B. }2\)  & \(\text{C. }3\)  & \(\text{D. }4\)
        \end{tabular*}\par}

        \item 在四棱锥 $P$-$ABCD$ 中，底面 $ABCD$ 为正方形，$AB=4$，$PC=PD=3$，$\angle PCA=45^\circ$，则 $\triangle PBC$ 的面积为
       
        {\centering\begin{tabular*}{\linewidth}{@{}@{\extracolsep{\fill}}llll@{}}
          \(\text{A. }2\sqrt{2}\)  & \(\text{B. }3\sqrt{2}\)  & \(\text{C. }4\sqrt{2}\)  & \(\text{D. }5\sqrt{2}\)
        \end{tabular*}\par}

        \item 已知椭圆 $\dfrac{x^2}{9}+\dfrac{y^2}{6}=1$，$F_1,F_2$ 为两个焦点，$O$ 为原点，$P$ 为椭圆上一点，$\cos\angle F_1PF_2=\dfrac{3}{5}$，则 $|PO|=$
       
        {\centering\begin{tabular*}{\linewidth}{@{}@{\extracolsep{\fill}}llll@{}}
          \(\text{A. }\dfrac{2}{5}\)  & \(\text{B. }\dfrac{\sqrt{30}}{2}\)  & \(\text{C. }\dfrac{3}{5}\)  & \(\text{D. }\dfrac{\sqrt{35}}{2}\)
        \end{tabular*}\par}
    \end{enumerate}
\end{description}

\begin{description}
    \item[二、填空题] 本题共4小题，每小题5分，共20分．
    \begin{enumerate}[leftmargin=5pt]
      \addtocounter{enumi}{+12}
        \item 若 $y=(x-1)^2+ax+\sin\bigl(x+\frac{\pi}{2}\bigr)$ 为偶函数，则 $a=$ \rule[-0.3ex]{3em}{0.4pt}．
        \item 设 $x,y$ 满足约束条件 $\begin{cases}-2x+3y\leqslant3,\\3x-2y\leqslant3,\\x+y\geqslant1,\end{cases}$ 则 $z=3x+2y$ 的最大值为 \rule[-0.3ex]{3em}{0.4pt}．
        \item 在正方体 $ABCD$-$A_1B_1C_1D_1$ 中，$E,F$ 分别为 $AB,C_1D_1$ 的中点．以 $EF$ 为直径的球的球面与该正方体的棱共有 \rule[-0.3ex]{3em}{0.4pt} 个公共点．
        \item 在 $\triangle ABC$ 中，$AB=2$，$\angle BAC=60^\circ$，$BC=\sqrt{6}$，$D$ 为 $BC$ 上一点，$AD$ 为 $\angle BAC$ 的角平分线，则 $AD=$ \rule[-0.3ex]{3em}{0.4pt}．
    \end{enumerate}
\end{description}

\begin{description}
    \item[三、解答题] 本题共7小题，共70分．解答应写出文字说明、证明过程或演算步骤．第17\textasciitilde21题为必考题，第22、23题为选考题．
    \begin{enumerate}[leftmargin=5pt]
      \addtocounter{enumi}{+16}
        \item （12分）

        已知数列 $\{a_n\}$ 中，$a_2=1$，设 $S_n$ 为 $\{a_n\}$ 前 $n$ 项和，$2S_n=na_n$．

        \begin{enumerate}[label=(\arabic*)]
            \item 求 $\{a_n\}$ 的通项公式；
            \item 求数列 $\bigl\{\frac{a_n+1}{2^n}\bigr\}$ 的前 $n$ 项和 $T_n$．
        \end{enumerate}

        \item （12分）

        如图，在三棱柱 $ABC$-$A_1B_1C_1$ 中，$A_1C\perp$ 平面 $ABC$，$\angle ACB=90^\circ$，$AA_1=2$，$A_1$ 到平面 $BCC_1B_1$ 的距离为 $1$．

        \textsf{[插入图片：三棱柱示意图]}

        \begin{enumerate}[label=(\arabic*)]
            \item 证明：$A_1C=AC$；
            \item 已知 $AA_1$ 与 $BB_1$ 距离为 $2$，求 $AB_1$ 与平面 $BCC_1B_1$ 所成角的正弦值．
        \end{enumerate}

        \item （12分）

        为探究某药物对小鼠的生长抑制作用，将40只小鼠均分为两组，分别为对照组（不加药物）和实验组（加药物）．

        \textsf{[插入图片：数据表格和直方图]}

        \begin{enumerate}[label=(\arabic*)]
            \item 设对照组、实验组小鼠体重增加量的中位数分别为 $m_0,m_1$，求 $P(m_0\geqslant m_1)$；
            \item 当 $p$ 取何值时，有 $95\%$ 的把握认为药物对小鼠生长有抑制作用．
        \end{enumerate}

        \item （12分）

        已知直线 $x-2y+1=0$ 与抛物线 $C:y^2=2px\,(p>0)$ 交于 $A,B$ 两点，$|AB|=4\sqrt{15}$．

        \begin{enumerate}[label=(\arabic*)]
            \item 求 $p$；
            \item 设 $F$ 为 $C$ 的焦点，$M,N$ 为 $C$ 上两点，且 $\overrightarrow{FM}\cdot\overrightarrow{FN}=0$，求 $\triangle MFN$ 面积的最小值．
        \end{enumerate}

        \item （12分）

        已知 $f(x)=ax-\dfrac{\sin x}{\cos^3 x}$，$x\in\bigl(0,\frac{\pi}{2}\bigr)$．

        \begin{enumerate}[label=(\arabic*)]
            \item 当 $a=8$ 时，讨论 $f(x)$ 的单调性；
            \item 若 $f(x)<\sin2x$，求 $a$ 的取值范围．
        \end{enumerate}

        \item （10分）[选修4-4：坐标系与参数方程]

        已知点 $P(2,1)$，直线 $l:\begin{cases}x=2+t\cos\alpha,\\y=1+t\sin\alpha\end{cases}$（$t$ 为参数），$\alpha$ 为 $l$ 的倾斜角．$l$ 与 $x$ 轴正半轴、$y$ 轴正半轴分别交于 $A,B$，$|PA|\cdot|PB|=4$．

        \begin{enumerate}[label=(\arabic*)]
            \item 求 $\alpha$；
            \item 以坐标原点为极点，$x$ 轴正半轴为极轴建立极坐标系，求 $l$ 的极坐标方程．
        \end{enumerate}

        \item （10分）[选修4-5：不等式选讲]

        已知 $f(x)=2|x-a|-a$，$a>0$．

        \begin{enumerate}[label=(\arabic*)]
            \item 求不等式 $f(x)<x$ 的解集；
            \item 若曲线 $y=f(x)$ 与 $x$ 轴所围成的图形的面积为 $2$，求 $a$．
        \end{enumerate}
    \end{enumerate}
\end{description}
\clearpage

\subsection{2023年全国甲卷（文）}\label{gk:2023年全国甲卷（文）}

\begin{description}
    \item[一、选择题] 本题共12小题，每小题5分，共60分．在每小题给出的四个选项中，只有一项是符合题目要求的．
    \begin{enumerate}[leftmargin=5pt]
        \item 设全集 $U=\{1,2,3,4,5\}$，集合 $M=\{1,4\}$，$N=\{2,5\}$，则 $N\cup\complement_U M=$
       
        {\centering\begin{tabular*}{\linewidth}{@{}@{\extracolsep{\fill}}llll@{}}
          \(\text{A. }\{2,3,5\}\)  & \(\text{B. }\{1,3,4\}\)  & \(\text{C. }\{1,2,4,5\}\)  & \(\text{D. }\{2,3,4,5\}\)
        \end{tabular*}\par}

        \item $\dfrac{5(1+\i^3)}{(2+\i)(2-\i)}=$
       
        {\centering\begin{tabular*}{\linewidth}{@{}@{\extracolsep{\fill}}llll@{}}
          \(\text{A. }-1\)  & \(\text{B. }1\)  & \(\text{C. }1-\i\)  & \(\text{D. }1+\i\)
        \end{tabular*}\par}

        \item 已知向量 $\bm a=(3,1)$，$\bm b=(2,2)$，则 $\cos\langle\bm a+\bm b,\bm a-\bm b\rangle=$
       
        {\centering\begin{tabular*}{\linewidth}{@{}@{\extracolsep{\fill}}llll@{}}
          \(\text{A. }\dfrac{1}{17}\)  & \(\text{B. }\dfrac{\sqrt{17}}{17}\)  & \(\text{C. }\dfrac{\sqrt{5}}{5}\)  & \(\text{D. }\dfrac{2\sqrt{5}}{5}\)
        \end{tabular*}\par}

        \item 某校为举办活动，要求在校门口按男女相间的方式站成两列．每列均为 $n\,(n\in\mathbb{N}^*)$ 人，则不同的站法种数为
       
        {\centering\begin{tabular*}{\linewidth}{@{}@{\extracolsep{\fill}}llll@{}}
          \(\text{A. }2(n!)^2\)  & \(\text{B. }n!(n+1)!\)  & \(\text{C. }(n!)^2\)  & \(\text{D. }2\cdot n!\)
        \end{tabular*}\par}

        \item 记 $S_n$ 为等差数列 $\{a_n\}$ 的前 $n$ 项和．若 $a_2+a_6=10$，$a_4a_8=45$，则 $S_5=$
       
        {\centering\begin{tabular*}{\linewidth}{@{}@{\extracolsep{\fill}}llll@{}}
          \(\text{A. }25\)  & \(\text{B. }22\)  & \(\text{C. }20\)  & \(\text{D. }15\)
        \end{tabular*}\par}

        \item 执行下面的程序框图，输出的 $B=$
        
        \textsf{[插入图片：程序框图]}

        {\centering\begin{tabular*}{\linewidth}{@{}@{\extracolsep{\fill}}llll@{}}
          \(\text{A. }21\)  & \(\text{B. }34\)  & \(\text{C. }55\)  & \(\text{D. }89\)
        \end{tabular*}\par}

        \item 设 $F_1,F_2$ 为椭圆 $C:\dfrac{x^2}{5}+y^2=1$ 的两个焦点，点 $P$ 在 $C$ 上，若 $\overrightarrow{PF_1}\cdot\overrightarrow{PF_2}=0$，则 $|PF_1|\cdot|PF_2|=$
       
        {\centering\begin{tabular*}{\linewidth}{@{}@{\extracolsep{\fill}}llll@{}}
          \(\text{A. }1\)  & \(\text{B. }2\)  & \(\text{C. }4\)  & \(\text{D. }5\)
        \end{tabular*}\par}

        \item 曲线 $y=\dfrac{\e^x}{x+1}$ 在点 $\bigl(1,\frac{\e}{2}\bigr)$ 处的切线方程为
       
        {\centering\begin{tabular*}{\linewidth}{@{}@{\extracolsep{\fill}}llll@{}}
          \(\text{A. }y=\dfrac{\e}{4}x\)  & \(\text{B. }y=\dfrac{\e}{2}x\)  & \(\text{C. }y=\dfrac{\e}{4}x+\dfrac{\e}{4}\)  & \(\text{D. }y=\dfrac{\e}{2}x+\dfrac{3\e}{4}\)
        \end{tabular*}\par}

        \item 已知双曲线 $\dfrac{x^2}{a^2}-\dfrac{y^2}{b^2}=1\,(a>0,b>0)$ 的离心率为 $\sqrt{5}$，其中一条渐近线与圆 $(x-2)^2+(y-3)^2=1$ 交于 $A,B$ 两点，则 $|AB|=$
       
        {\centering\begin{tabular*}{\linewidth}{@{}@{\extracolsep{\fill}}llll@{}}
          \(\text{A. }\dfrac{1}{5}\)  & \(\text{B. }\dfrac{\sqrt{5}}{5}\)  & \(\text{C. }\dfrac{2\sqrt{5}}{5}\)  & \(\text{D. }\dfrac{4\sqrt{5}}{5}\)
        \end{tabular*}\par}

        \item 在三棱锥 $P$-$ABC$ 中，$\triangle ABC$ 是边长为 $2$ 的等边三角形，$PA=PB=2$，$PC=\sqrt{6}$，则该棱锥的体积为
       
        {\centering\begin{tabular*}{\linewidth}{@{}@{\extracolsep{\fill}}llll@{}}
          \(\text{A. }1\)  & \(\text{B. }\sqrt{3}\)  & \(\text{C. }2\)  & \(\text{D. }3\)
        \end{tabular*}\par}

        \item 已知函数 $f(x)=\e^{-(x-1)^2}$．记 $a=f\bigl(\frac{\sqrt{2}}{2}\bigr)$，$b=f\bigl(\frac{\sqrt{3}}{2}\bigr)$，$c=f\bigl(\frac{\sqrt{6}}{2}\bigr)$，则
       
        {\centering\begin{tabular*}{\linewidth}{@{}@{\extracolsep{\fill}}llll@{}}
          \(\text{A. }b>c>a\)  & \(\text{B. }b>a>c\)  & \(\text{C. }c>a>b\)  & \(\text{D. }c>b>a\)
        \end{tabular*}\par}

        \item 函数 $y=f(x)$ 的图象由 $y=\cos\bigl(2x+\frac{\pi}{6}\bigr)$ 的图象向左平移 $\frac{\pi}{6}$ 个单位长度得到，则 $y=f(x)$ 的图象与直线 $y=\frac{1}{2}x-\frac{1}{2}$ 的交点个数为
       
        {\centering\begin{tabular*}{\linewidth}{@{}@{\extracolsep{\fill}}llll@{}}
          \(\text{A. }1\)  & \(\text{B. }2\)  & \(\text{C. }3\)  & \(\text{D. }4\)
        \end{tabular*}\par}
    \end{enumerate}
\end{description}

\begin{description}
    \item[二、填空题] 本题共4小题，每小题5分，共20分．
    \begin{enumerate}[leftmargin=5pt]
      \addtocounter{enumi}{+12}
        \item 记 $S_n$ 为等比数列 $\{a_n\}$ 的前 $n$ 项和．若 $8S_6=7S_3$，则 $\{a_n\}$ 的公比为 \rule[-0.3ex]{3em}{0.4pt}．
        \item 若 $y=(x-1)^2+ax+\sin\bigl(x+\frac{\pi}{2}\bigr)$ 为偶函数，则 $a=$ \rule[-0.3ex]{3em}{0.4pt}．
        \item 若 $x,y$ 满足约束条件 $\begin{cases}-2x+3y\leqslant3,\\3x-2y\leqslant3,\\x+y\geqslant1,\end{cases}$ 则 $z=3x+2y$ 的最大值为 \rule[-0.3ex]{3em}{0.4pt}．
        \item 在正方体 $ABCD$-$A_1B_1C_1D_1$ 中，$AB=4$，$O$ 为 $A_1C_1$ 的中点．若该正方体的棱与球 $O$ 的球面有公共点，则球 $O$ 的半径的取值范围是 \rule[-0.3ex]{3em}{0.4pt}．
    \end{enumerate}
\end{description}

\begin{description}
    \item[三、解答题] 本题共7小题，共70分．解答应写出文字说明、证明过程或演算步骤．第17\textasciitilde21题为必考题，第22、23题为选考题．
    \begin{enumerate}[leftmargin=5pt]
      \addtocounter{enumi}{+16}
        \item （12分）

        记 $\triangle ABC$ 的内角 $A,B,C$ 的对边分别为 $a,b,c$，已知 $\dfrac{b^2+c^2-a^2}{\cos A}=2$．

        \begin{enumerate}[label=(\arabic*)]
            \item 求 $bc$；
            \item 若 $\dfrac{a\cos B-b\cos A}{a\cos B+b\cos A}-\dfrac{b}{c}=1$，求 $\triangle ABC$ 面积．
        \end{enumerate}

        \item （12分）

        如图，在三棱柱 $ABC$-$A_1B_1C_1$ 中，$A_1C\perp$ 平面 $ABC$，$\angle ACB=90^\circ$．

        \textsf{[插入图片：三棱柱示意图]}

        \begin{enumerate}[label=(\arabic*)]
            \item 证明：平面 $ACC_1A_1\perp$ 平面 $BB_1C_1C$；
            \item 设 $AB=A_1B$，$AA_1=2$，求四棱锥 $A_1$-$BB_1C_1C$ 的高．
        \end{enumerate}

        \item （12分）

        一项试验旨在研究臭氧效应，试验方案如下：选40只小白鼠，随机地将其中20只分配到试验组，另外20只分配到对照组，试验组…

        \textsf{[插入图片：数据表格]}

        \begin{enumerate}[label=(\arabic*)]
            \item 计算对照组和试验组的平均体重，判断臭氧对小白鼠体重的影响；
            \item 根据数据，能否有 $95\%$ 的把握认为臭氧对小白鼠体重有影响？
        \end{enumerate}

        \item （12分）

        已知函数 $f(x)=ax-\dfrac{\sin x}{\cos^2 x}$，$x\in\bigl(0,\frac{\pi}{2}\bigr)$．

        \begin{enumerate}[label=(\arabic*)]
            \item 当 $a=1$ 时，讨论 $f(x)$ 的单调性；
            \item 若 $f(x)+\sin x<0$，求 $a$ 的取值范围．
        \end{enumerate}

        \item （12分）

        已知直线 $x-2y+1=0$ 与抛物线 $C:y^2=2px\,(p>0)$ 交于 $A,B$ 两点，$|AB|=4\sqrt{15}$．

        \begin{enumerate}[label=(\arabic*)]
            \item 求 $p$；
            \item 设 $F$ 为 $C$ 的焦点，$M,N$ 为 $C$ 上两点，且 $\overrightarrow{FM}\cdot\overrightarrow{FN}=0$，求 $\triangle MFN$ 面积的最小值．
        \end{enumerate}

        \item （10分）[选修4-4：坐标系与参数方程]

        已知点 $P(2,1)$，直线 $l:\begin{cases}x=2+t\cos\alpha,\\y=1+t\sin\alpha\end{cases}$（$t$ 为参数），$\alpha$ 为 $l$ 的倾斜角．$l$ 与 $x$ 轴正半轴、$y$ 轴正半轴分别交于 $A,B$，$|PA|\cdot|PB|=4$．

        \begin{enumerate}[label=(\arabic*)]
            \item 求 $\alpha$；
            \item 以坐标原点为极点，$x$ 轴正半轴为极轴建立极坐标系，求 $l$ 的极坐标方程．
        \end{enumerate}

        \item （10分）[选修4-5：不等式选讲]

        设 $a>0$，函数 $f(x)=2|x-a|-a$．

        \begin{enumerate}[label=(\arabic*)]
            \item 求不等式 $f(x)<x$ 的解集；
            \item 若曲线 $y=f(x)$ 与 $x$ 轴所围成的图形的面积为 $2$，求 $a$．
        \end{enumerate}
    \end{enumerate}
\end{description}
\clearpage

\subsection{2023年全国乙卷（理）}\label{gk:2023年全国乙卷（理）}

\begin{description}
    \item[一、选择题] 本题共12小题，每小题5分，共60分．在每小题给出的四个选项中，只有一项是符合题目要求的．
    \begin{enumerate}[leftmargin=5pt]
        \item 设 $z=\dfrac{2+\i}{1+\i^2+\i^5}$，则 $\bar{z}=$
       
        {\centering\begin{tabular*}{\linewidth}{@{}@{\extracolsep{\fill}}llll@{}}
          \(\text{A. }1-2\i\)  & \(\text{B. }1+2\i\)  & \(\text{C. }2-\i\)  & \(\text{D. }2+\i\)
        \end{tabular*}\par}

        \item 设集合 $U=\mathbb{R}$，集合 $M=\{x\mid x<1\}$，$N=\{x\mid -1<x<2\}$，则 $\{x\mid x\geqslant2\}=$
       
        {\centering\begin{tabular*}{\linewidth}{@{}@{\extracolsep{\fill}}llll@{}}
          \(\text{A. }\complement_U(M\cup N)\)  & \(\text{B. }N\cup\complement_U M\)  & \(\text{C. }\complement_U(M\cap N)\)  & \(\text{D. }M\cup\complement_U N\)
        \end{tabular*}\par}

        \item 如图，网格纸上绘制的是一个零件的三视图，网格小正方形的边长为1，则该零件的表面积为
        
        \textsf{[插入图片：三视图]}

        {\centering\begin{tabular*}{\linewidth}{@{}@{\extracolsep{\fill}}llll@{}}
          \(\text{A. }24\)  & \(\text{B. }26\)  & \(\text{C. }28\)  & \(\text{D. }30\)
        \end{tabular*}\par}

        \item 已知 $f(x)=\dfrac{x\e^x}{\e^{ax}-1}$ 是偶函数，则 $a=$
       
        {\centering\begin{tabular*}{\linewidth}{@{}@{\extracolsep{\fill}}llll@{}}
          \(\text{A. }-2\)  & \(\text{B. }-1\)  & \(\text{C. }1\)  & \(\text{D. }2\)
        \end{tabular*}\par}

        \item 设 $O$ 为平面坐标系的坐标原点，在区域 $\{(x,y)\mid 1\leqslant x^2+y^2\leqslant4\}$ 内随机取一点，记该点为 $A$，则直线 $OA$ 的倾斜角不大于 $\frac{\pi}{4}$ 的概率为
       
        {\centering\begin{tabular*}{\linewidth}{@{}@{\extracolsep{\fill}}llll@{}}
          \(\text{A. }\dfrac{1}{8}\)  & \(\text{B. }\dfrac{1}{6}\)  & \(\text{C. }\dfrac{1}{4}\)  & \(\text{D. }\dfrac{1}{2}\)
        \end{tabular*}\par}

        \item 已知函数 $f(x)=\sin(\omega x+\varphi)$ 在区间 $\bigl(\frac{\pi}{6},\frac{2\pi}{3}\bigr)$ 单调递增，直线 $x=\frac{\pi}{6}$ 和 $x=\frac{2\pi}{3}$ 为函数 $y=f(x)$ 的图像的两条对称轴，则 $f\bigl(-\frac{5\pi}{12}\bigr)=$
       
        {\centering\begin{tabular*}{\linewidth}{@{}@{\extracolsep{\fill}}llll@{}}
          \(\text{A. }-\dfrac{\sqrt{3}}{2}\)  & \(\text{B. }-\dfrac{1}{2}\)  & \(\text{C. }\dfrac{1}{2}\)  & \(\text{D. }\dfrac{\sqrt{3}}{2}\)
        \end{tabular*}\par}

        \item 甲乙两位同学从6种课外读物中各自选读2种，则这两人选读的课外读物中恰有一种相同的选法共有
       
        {\centering\begin{tabular*}{\linewidth}{@{}@{\extracolsep{\fill}}llll@{}}
          \(\text{A. }30\) 种  & \(\text{B. }60\) 种  & \(\text{C. }120\) 种  & \(\text{D. }240\) 种
        \end{tabular*}\par}

        \item 已知圆锥 $PO$ 的底面半径为 $\sqrt{3}$，$O$ 为底面圆心，$PA,PB$ 为圆锥的母线，$\angle AOB=120^\circ$，若 $\triangle PAB$ 的面积为 $\dfrac{9\sqrt{3}}{4}$，则该圆锥的体积为
       
        {\centering\begin{tabular*}{\linewidth}{@{}@{\extracolsep{\fill}}llll@{}}
          \(\text{A. }\pi\)  & \(\text{B. }\sqrt{6}\pi\)  & \(\text{C. }3\pi\)  & \(\text{D. }3\sqrt{6}\pi\)
        \end{tabular*}\par}

        \item 已知 $\triangle ABC$ 为等腰直角三角形，$AB$ 为斜边，$\triangle ABD$ 为等边三角形，若二面角 $C$-$AB$-$D$ 为 $150^\circ$，则直线 $CD$ 与平面 $ABC$ 所成角的正切值为
       
        {\centering\begin{tabular*}{\linewidth}{@{}@{\extracolsep{\fill}}llll@{}}
          \(\text{A. }\dfrac{1}{5}\)  & \(\text{B. }\dfrac{\sqrt{2}}{5}\)  & \(\text{C. }\dfrac{\sqrt{3}}{5}\)  & \(\text{D. }\dfrac{2}{5}\)
        \end{tabular*}\par}

        \item 已知等差数列 $\{a_n\}$ 的公差为 $\dfrac{2\pi}{3}$，集合 $S=\{\cos a_n\mid n\in\mathbb{N}^*\}$，若 $S=\{a,b\}$，则 $ab=$
       
        {\centering\begin{tabular*}{\linewidth}{@{}@{\extracolsep{\fill}}llll@{}}
          \(\text{A. }-1\)  & \(\text{B. }-\dfrac{1}{2}\)  & \(\text{C. }0\)  & \(\text{D. }\dfrac{1}{2}\)
        \end{tabular*}\par}

        \item 设 $A,B$ 为双曲线 $x^2-\dfrac{y^2}{9}=1$ 上两点，下列四个点中，可为线段 $AB$ 中点的是
       
        {\centering\begin{tabular*}{\linewidth}{@{}@{\extracolsep{\fill}}llll@{}}
          \(\text{A. }(1,1)\)  & \(\text{B. }(-1,2)\)  & \(\text{C. }(1,3)\)  & \(\text{D. }(-1,-4)\)
        \end{tabular*}\par}

        \item 已知 $\odot O$ 的半径为 $1$，直线 $PA$ 与 $\odot O$ 相切于点 $A$，直线 $PB$ 与 $\odot O$ 交于 $B,C$ 两点，$D$ 为 $BC$ 中点．若 $|PO|=\sqrt{2}$，则 $\overrightarrow{PA}\cdot\overrightarrow{PD}$ 的最大值为
       
        {\centering\begin{tabular*}{\linewidth}{@{}@{\extracolsep{\fill}}llll@{}}
          \(\text{A. }\dfrac{1+\sqrt{2}}{2}\)  & \(\text{B. }\dfrac{1+2\sqrt{2}}{2}\)  & \(\text{C. }1+\sqrt{2}\)  & \(\text{D. }2+\sqrt{2}\)
        \end{tabular*}\par}
    \end{enumerate}
\end{description}

\begin{description}
    \item[二、填空题] 本题共4小题，每小题5分，共20分．
    \begin{enumerate}[leftmargin=5pt]
      \addtocounter{enumi}{+12}
        \item 已知点 $A(1,\sqrt{5})$ 在抛物线 $C:y^2=2px$ 上，则 $A$ 到 $C$ 的准线的距离为 \rule[-0.3ex]{3em}{0.4pt}．
        \item 若 $x,y$ 满足约束条件 $\begin{cases}x-3y\leqslant-7,\\x+2y\leqslant9,\\3x+y\geqslant7,\end{cases}$ 则 $z=2x-y$ 的最大值为 \rule[-0.3ex]{3em}{0.4pt}．
        \item 已知 $\{a_n\}$ 为等比数列，$a_2a_4a_5=a_3a_6$，$a_9a_{10}=-8$，则 $a_7=$ \rule[-0.3ex]{3em}{0.4pt}．
        \item 设 $a\in(0,1)$，若函数 $f(x)=a^x+(1+a)^x$ 在 $(0,+\infty)$ 上单调递增，则 $a$ 的取值范围是 \rule[-0.3ex]{3em}{0.4pt}．
    \end{enumerate}
\end{description}

\begin{description}
    \item[三、解答题] 本题共7小题，共70分．解答应写出文字说明、证明过程或演算步骤．第17\textasciitilde21题为必考题，第22、23题为选考题．
    \begin{enumerate}[leftmargin=5pt]
      \addtocounter{enumi}{+16}
        \item （12分）

        某厂为比较甲乙两种工艺对橡胶产品伸缩率的处理效应，进行10次配对试验，每次配对试验选用材质相同的两个橡胶产品，随机地选其中一个用甲工艺处理，另一个用乙工艺处理…

        \textsf{[插入图片：试验数据表]}

        \begin{enumerate}[label=(\arabic*)]
            \item 求试验组和对照组的伸缩率，判断哪种工艺更好；
            \item 计算并比较两种工艺处理后的橡胶产品伸缩率的方差．
        \end{enumerate}

        \item （12分）

        在 $\triangle ABC$ 中，已知 $\angle BAC=120^\circ$，$AB=2$，$AC=1$．

        \begin{enumerate}[label=(\arabic*)]
            \item 求 $\sin\angle ABC$；
            \item 若 $D$ 为 $BC$ 上一点，且 $\angle BAD=90^\circ$，求 $\triangle ADC$ 的面积．
        \end{enumerate}

        \item （12分）

        如图，在三棱锥 $P$-$ABC$ 中，$AB\perp BC$，$AB=2$，$BC=2\sqrt{2}$，$PB=PC=\sqrt{6}$，$BP,AP,BC$ 的中点分别为 $D,E,O$，$AD=\sqrt{5}DO$，点 $F$ 在 $AC$ 上，$BF\perp AO$．

        \textsf{[插入图片：三棱锥示意图]}

        \begin{enumerate}[label=(\arabic*)]
            \item 证明：$EF\parallel$ 平面 $ADO$；
            \item 证明：平面 $ADO\perp$ 平面 $BEF$；
            \item 求二面角 $D$-$AO$-$C$ 的正弦值．
        \end{enumerate}

        \item （12分）

        已知椭圆 $C:\dfrac{y^2}{a^2}+\dfrac{x^2}{b^2}=1\,(a>b>0)$ 的离心率为 $\dfrac{\sqrt{5}}{3}$，点 $A(-2,0)$ 在 $C$ 上．

        \begin{enumerate}[label=(\arabic*)]
            \item 求 $C$ 的方程；
            \item 过点 $(-2,3)$ 的直线交 $C$ 于 $P,Q$ 两点，直线 $AP,AQ$ 与 $y$ 轴的交点分别为 $M,N$，证明：线段 $MN$ 的中点为定点．
        \end{enumerate}

        \item （12分）

        已知函数 $f(x)=\bigl(\frac{1}{x}+a\bigr)\ln(1+x)$．

        \begin{enumerate}[label=(\arabic*)]
            \item 当 $a=-1$ 时，求曲线 $y=f(x)$ 在点 $(1,f(1))$ 处的切线方程；
            \item 是否存在 $a,b$，使得曲线 $y=f\bigl(\frac{1}{x}\bigr)$ 关于直线 $x=b$ 对称？若存在，求 $a,b$；若不存在，说明理由；
            \item 若 $f(x)$ 在 $(0,+\infty)$ 存在极值，求 $a$ 的取值范围．
        \end{enumerate}

        \item （10分）[选修4-4：坐标系与参数方程]

        在直角坐标系 $xOy$ 中，曲线 $C_1$ 的参数方程为 $\begin{cases}x=\dfrac{2\cos\theta}{\sin^2\theta},\\y=\dfrac{2}{\sin\theta}\end{cases}$（$\theta$ 为参数，$\frac{\pi}{2}<\theta<\pi$ 或 $\pi<\theta<\frac{3\pi}{2}$）．以坐标原点为极点，$x$ 轴正半轴为极轴建立极坐标系，已知射线 $l:\theta=\alpha\,(\rho\geqslant0)$ 与 $C_1$ 交于 $A$，与 $C_2:\rho=2\sin\theta$ 交于 $B$，$|AB|=\sqrt{2}$．

        \begin{enumerate}[label=(\arabic*)]
            \item 求 $\alpha$；
            \item 求 $C_1$ 的极坐标方程．
        \end{enumerate}

        \item （10分）[选修4-5：不等式选讲]

        已知 $f(x)=2|x|+|x-2|$．

        \begin{enumerate}[label=(\arabic*)]
            \item 求不等式 $f(x)\leqslant6-x$ 的解集；
            \item 在直角坐标系 $xOy$ 中，求不等式组 $\begin{cases}f(x)\leqslant y,\\x+y-6\leqslant0\end{cases}$ 所确定的平面区域的面积．
        \end{enumerate}
    \end{enumerate}
\end{description}
\clearpage

\subsection{2023年全国乙卷（文）}\label{gk:2023年全国乙卷（文）}

\begin{description}
    \item[一、选择题] 本题共12小题，每小题5分，共60分．在每小题给出的四个选项中，只有一项是符合题目要求的．
    \begin{enumerate}[leftmargin=5pt]
        \item $|2+\i^2+2\i^3|=$
       
        {\centering\begin{tabular*}{\linewidth}{@{}@{\extracolsep{\fill}}llll@{}}
          \(\text{A. }1\)  & \(\text{B. }2\)  & \(\text{C. }\sqrt{5}\)  & \(\text{D. }5\)
        \end{tabular*}\par}

        \item 设全集 $U=\{0,1,2,4,6,8\}$，集合 $M=\{0,4,6\}$，$N=\{0,1,6\}$，则 $M\cup\complement_U N=$
       
        {\centering\begin{tabular*}{\linewidth}{@{}@{\extracolsep{\fill}}llll@{}}
          \(\text{A. }\{0,2,4,6,8\}\)  & \(\text{B. }\{0,1,4,6,8\}\)  & \(\text{C. }\{1,2,4,6,8\}\)  & \(\text{D. }U\)
        \end{tabular*}\par}

        \item 如图，网格纸上绘制的是一个零件的三视图，网格小正方形的边长为1，则该零件的表面积为
        
        \textsf{[插入图片：三视图]}

        {\centering\begin{tabular*}{\linewidth}{@{}@{\extracolsep{\fill}}llll@{}}
          \(\text{A. }24\)  & \(\text{B. }26\)  & \(\text{C. }28\)  & \(\text{D. }30\)
        \end{tabular*}\par}

        \item 在 $\triangle ABC$ 中，内角 $A,B,C$ 的对边分别是 $a,b,c$，若 $a\cos B-b\cos A=c$，且 $C=\dfrac{\pi}{5}$，则 $B=$
       
        {\centering\begin{tabular*}{\linewidth}{@{}@{\extracolsep{\fill}}llll@{}}
          \(\text{A. }\dfrac{\pi}{10}\)  & \(\text{B. }\dfrac{\pi}{5}\)  & \(\text{C. }\dfrac{3\pi}{10}\)  & \(\text{D. }\dfrac{2\pi}{5}\)
        \end{tabular*}\par}

        \item 已知 $f(x)=\dfrac{x\e^x}{\e^{ax}-1}$ 是偶函数，则 $a=$
       
        {\centering\begin{tabular*}{\linewidth}{@{}@{\extracolsep{\fill}}llll@{}}
          \(\text{A. }-2\)  & \(\text{B. }-1\)  & \(\text{C. }1\)  & \(\text{D. }2\)
        \end{tabular*}\par}

        \item 正方形 $ABCD$ 的边长是 $2$，$E$ 是 $AB$ 的中点，则 $\overrightarrow{EC}\cdot\overrightarrow{ED}=$
       
        {\centering\begin{tabular*}{\linewidth}{@{}@{\extracolsep{\fill}}llll@{}}
          \(\text{A. }\sqrt{5}\)  & \(\text{B. }3\)  & \(\text{C. }2\sqrt{5}\)  & \(\text{D. }5\)
        \end{tabular*}\par}

        \item 设 $O$ 为平面坐标系的坐标原点，在区域 $\{(x,y)\mid 1\leqslant x^2+y^2\leqslant4\}$ 内随机取一点 $A$，则直线 $OA$ 的倾斜角不大于 $\frac{\pi}{4}$ 的概率为
       
        {\centering\begin{tabular*}{\linewidth}{@{}@{\extracolsep{\fill}}llll@{}}
          \(\text{A. }\dfrac{1}{8}\)  & \(\text{B. }\dfrac{1}{6}\)  & \(\text{C. }\dfrac{1}{4}\)  & \(\text{D. }\dfrac{1}{2}\)
        \end{tabular*}\par}

        \item 函数 $f(x)=x^3+ax+2$ 存在3个零点，则 $a$ 的取值范围是
       
        {\centering\begin{tabular*}{\linewidth}{@{}@{\extracolsep{\fill}}llll@{}}
          \(\text{A. }(-\infty,-2)\)  & \(\text{B. }(-\infty,-3)\)  & \(\text{C. }(-4,-1)\)  & \(\text{D. }(-3,0)\)
        \end{tabular*}\par}

        \item 某学校举办作文比赛，共6个主题，每位参赛同学从中随机抽取一个主题准备作文．甲乙两位同学都参加比赛，则他们抽到的主题不同的概率为
       
        {\centering\begin{tabular*}{\linewidth}{@{}@{\extracolsep{\fill}}llll@{}}
          \(\text{A. }\dfrac{5}{6}\)  & \(\text{B. }\dfrac{2}{3}\)  & \(\text{C. }\dfrac{1}{2}\)  & \(\text{D. }\dfrac{1}{3}\)
        \end{tabular*}\par}

        \item 已知函数 $f(x)=\sin(\omega x+\varphi)$ 在区间 $\bigl(\frac{\pi}{6},\frac{2\pi}{3}\bigr)$ 单调递增，直线 $x=\frac{\pi}{6}$ 和 $x=\frac{2\pi}{3}$ 为函数 $y=f(x)$ 的图像的两条对称轴，则 $f\bigl(-\frac{5\pi}{12}\bigr)=$
       
        {\centering\begin{tabular*}{\linewidth}{@{}@{\extracolsep{\fill}}llll@{}}
          \(\text{A. }-\dfrac{\sqrt{3}}{2}\)  & \(\text{B. }-\dfrac{1}{2}\)  & \(\text{C. }\dfrac{1}{2}\)  & \(\text{D. }\dfrac{\sqrt{3}}{2}\)
        \end{tabular*}\par}

        \item 已知实数 $x,y$ 满足 $x^2+y^2-4x-2y-4=0$，则 $x-y$ 的最大值是
       
        {\centering\begin{tabular*}{\linewidth}{@{}@{\extracolsep{\fill}}llll@{}}
          \(\text{A. }1+\dfrac{3\sqrt{2}}{2}\)  & \(\text{B. }4\)  & \(\text{C. }1+3\sqrt{2}\)  & \(\text{D. }7\)
        \end{tabular*}\par}

        \item 设 $A,B$ 为双曲线 $x^2-\dfrac{y^2}{9}=1$ 上两点，下列四个点中，可为线段 $AB$ 中点的是
       
        {\centering\begin{tabular*}{\linewidth}{@{}@{\extracolsep{\fill}}llll@{}}
          \(\text{A. }(1,1)\)  & \(\text{B. }(-1,2)\)  & \(\text{C. }(1,3)\)  & \(\text{D. }(-1,-4)\)
        \end{tabular*}\par}
    \end{enumerate}
\end{description}

\begin{description}
    \item[二、填空题] 本题共4小题，每小题5分，共20分．
    \begin{enumerate}[leftmargin=5pt]
      \addtocounter{enumi}{+12}
        \item 已知点 $A(1,\sqrt{5})$ 在抛物线 $C:y^2=2px$ 上，则 $A$ 到 $C$ 的准线的距离为 \rule[-0.3ex]{3em}{0.4pt}．
        \item 若 $\theta\in\bigl(0,\frac{\pi}{2}\bigr)$，$\tan\theta=\dfrac{1}{3}$，则 $\sin\theta-\cos\theta=$ \rule[-0.3ex]{3em}{0.4pt}．
        \item 若 $x,y$ 满足约束条件 $\begin{cases}x-3y\leqslant-7,\\x+2y\leqslant9,\\3x+y\geqslant7,\end{cases}$ 则 $z=2x-y$ 的最大值为 \rule[-0.3ex]{3em}{0.4pt}．
        \item 已知点 $S,A,B,C$ 均在半径为 $2$ 的球面上，$\triangle ABC$ 是边长为 $3$ 的等边三角形，$SA\perp$ 平面 $ABC$，则 $SA=$ \rule[-0.3ex]{3em}{0.4pt}．
    \end{enumerate}
\end{description}

\begin{description}
    \item[三、解答题] 本题共7小题，共70分．解答应写出文字说明、证明过程或演算步骤．第17\textasciitilde21题为必考题，第22、23题为选考题．
    \begin{enumerate}[leftmargin=5pt]
      \addtocounter{enumi}{+16}
        \item （12分）

        某厂为比较甲乙两种工艺对橡胶产品伸缩率的处理效应，进行10次配对试验…

        \textsf{[插入图片：试验数据表]}

        \begin{enumerate}[label=(\arabic*)]
            \item 计算试验组和对照组的伸缩率均值，判断哪种工艺更好；
            \item 计算两种工艺处理后伸缩率的样本方差，比较其稳定性．
        \end{enumerate}

        \item （12分）

        记 $S_n$ 为等差数列 $\{a_n\}$ 的前 $n$ 项和，已知 $a_2=11$，$S_{10}=40$．

        \begin{enumerate}[label=(\arabic*)]
            \item 求 $\{a_n\}$ 的通项公式；
            \item 求数列 $\{|a_n|\}$ 的前 $n$ 项和 $T_n$．
        \end{enumerate}

        \item （12分）

        如图，在三棱锥 $P$-$ABC$ 中，$AB\perp BC$，$AB=2$，$BC=2\sqrt{2}$，$PB=PC=\sqrt{6}$，$BP,AP,BC$ 的中点分别为 $D,E,O$，$AD=\sqrt{5}DO$，点 $F$ 在 $AC$ 上，$BF\perp AO$．

        \textsf{[插入图片：三棱锥示意图]}

        \begin{enumerate}[label=(\arabic*)]
            \item 证明：$EF\parallel$ 平面 $ADO$；
            \item 证明：平面 $ADO\perp$ 平面 $BEF$；
            \item 求二面角 $D$-$AO$-$C$ 的正弦值．
        \end{enumerate}

        \item （12分）

        已知函数 $f(x)=\bigl(\frac{1}{x}+a\bigr)\ln(1+x)$．

        \begin{enumerate}[label=(\arabic*)]
            \item 当 $a=-1$ 时，求曲线 $y=f(x)$ 在点 $(1,f(1))$ 处的切线方程；
            \item 若函数 $f(x)$ 在 $(0,+\infty)$ 单调递增，求 $a$ 的取值范围．
        \end{enumerate}

        \item （12分）

        已知椭圆 $C:\dfrac{y^2}{a^2}+\dfrac{x^2}{b^2}=1\,(a>b>0)$ 的离心率为 $\dfrac{\sqrt{5}}{3}$，点 $A(-2,0)$ 在 $C$ 上．

        \begin{enumerate}[label=(\arabic*)]
            \item 求 $C$ 的方程；
            \item 过点 $(-2,3)$ 的直线交 $C$ 于 $P,Q$ 两点，直线 $AP,AQ$ 与 $y$ 轴的交点分别为 $M,N$，证明：线段 $MN$ 的中点为定点．
        \end{enumerate}

        \item （10分）[选修4-4：坐标系与参数方程]

        在直角坐标系 $xOy$ 中，曲线 $C_1$ 的参数方程为 $\begin{cases}x=\dfrac{2\cos\theta}{\sin^2\theta},\\y=\dfrac{2}{\sin\theta}\end{cases}$（$\theta$ 为参数）．以坐标原点为极点，$x$ 轴正半轴为极轴建立极坐标系，已知射线 $l:\theta=\alpha$ 与 $C_1$ 交于 $A$，与 $C_2:\rho=2\sin\theta$ 交于 $B$，$|AB|=\sqrt{2}$．

        \begin{enumerate}[label=(\arabic*)]
            \item 求 $\alpha$；
            \item 求 $C_1$ 的极坐标方程．
        \end{enumerate}

        \item （10分）[选修4-5：不等式选讲]

        已知 $f(x)=2|x|+|x-2|$．

        \begin{enumerate}[label=(\arabic*)]
            \item 求不等式 $f(x)\leqslant6-x$ 的解集；
            \item 求不等式组 $\begin{cases}f(x)\leqslant y,\\x+y-6\leqslant0\end{cases}$ 所确定的平面区域的面积．
        \end{enumerate}
    \end{enumerate}
\end{description}
\clearpage

\subsection{2023年新高考I卷}\label{gk:2023年新高考I卷}

\begin{description}
    \item[一、选择题] 本题共8小题，每小题5分，共40分．在每小题给出的四个选项中，只有一项是符合题目要求的．
    \begin{enumerate}[leftmargin=5pt]
        \item 已知集合 $M=\{-2,-1,0,1,2\}$，$N=\{x\mid x^2-x-6\geqslant0\}$，则 $M\cap N=$
       
        {\centering\begin{tabular*}{\linewidth}{@{}@{\extracolsep{\fill}}llll@{}}
          \(\text{A. }\{-2,-1,0,1\}\)  & \(\text{B. }\{0,1,2\}\)  & \(\text{C. }\{-2\}\)  & \(\text{D. }\{2\}\)
        \end{tabular*}\par}

        \item 已知 $z=\dfrac{1-\i}{2+2\i}$，则 $z-\bar{z}=$
       
        {\centering\begin{tabular*}{\linewidth}{@{}@{\extracolsep{\fill}}llll@{}}
          \(\text{A. }-\i\)  & \(\text{B. }\i\)  & \(\text{C. }0\)  & \(\text{D. }1\)
        \end{tabular*}\par}

        \item 已知向量 $\bm a=(1,1)$，$\bm b=(1,-1)$．若 $(\bm a+\lambda\bm b)\perp(\bm a+\mu\bm b)$，则
       
        {\centering\begin{tabular*}{\linewidth}{@{}@{\extracolsep{\fill}}llll@{}}
          \(\text{A. }\lambda+\mu=1\)  & \(\text{B. }\lambda+\mu=-1\)  & \(\text{C. }\lambda\mu=1\)  & \(\text{D. }\lambda\mu=-1\)
        \end{tabular*}\par}

        \item 设函数 $f(x)=2^{x(x-a)}$ 在区间 $(0,1)$ 单调递减，则 $a$ 的取值范围是
       
        {\centering\begin{tabular*}{\linewidth}{@{}@{\extracolsep{\fill}}llll@{}}
          \(\text{A. }(-\infty,-2]\)  & \(\text{B. }[-2,0)\)  & \(\text{C. }(0,2]\)  & \(\text{D. }[2,+\infty)\)
        \end{tabular*}\par}

        \item 设椭圆 $C_1:\dfrac{x^2}{a^2}+y^2=1\,(a>1)$，$C_2:\dfrac{x^2}{4}+y^2=1$ 的离心率分别为 $e_1,e_2$．若 $e_2=\sqrt{3}e_1$，则 $a=$
       
        {\centering\begin{tabular*}{\linewidth}{@{}@{\extracolsep{\fill}}llll@{}}
          \(\text{A. }\dfrac{2\sqrt{3}}{3}\)  & \(\text{B. }\sqrt{2}\)  & \(\text{C. }\sqrt{3}\)  & \(\text{D. }\sqrt{6}\)
        \end{tabular*}\par}

        \item 过点 $(0,-2)$ 与圆 $x^2+y^2-4x-1=0$ 相切的两条直线的夹角为 $\alpha$，则 $\sin\alpha=$
       
        {\centering\begin{tabular*}{\linewidth}{@{}@{\extracolsep{\fill}}llll@{}}
          \(\text{A. }1\)  & \(\text{B. }\dfrac{\sqrt{15}}{4}\)  & \(\text{C. }\dfrac{\sqrt{10}}{4}\)  & \(\text{D. }\dfrac{\sqrt{6}}{4}\)
        \end{tabular*}\par}

        \item 记 $S_n$ 为数列 $\{a_n\}$ 的前 $n$ 项和，设甲：$\{a_n\}$ 为等差数列；乙：$\bigl\{\frac{S_n}{n}\bigr\}$ 为等差数列，则
       
        {\centering\begin{tabular*}{\linewidth}{@{}@{\extracolsep{\fill}}llll@{}}
          \(\text{A. }\)甲是乙的充分条件但不是必要条件  & \(\text{B. }\)甲是乙的必要条件但不是充分条件 \\
          \(\text{C. }\)甲是乙的充要条件  & \(\text{D. }\)甲既不是乙的充分条件也不是乙的必要条件
        \end{tabular*}\par}

        \item 已知 $\sin(\alpha-\beta)=\dfrac{1}{3}$，$\cos\alpha\sin\beta=\dfrac{1}{6}$，则 $\cos(2\alpha+2\beta)=$
       
        {\centering\begin{tabular*}{\linewidth}{@{}@{\extracolsep{\fill}}llll@{}}
          \(\text{A. }\dfrac{7}{9}\)  & \(\text{B. }\dfrac{1}{9}\)  & \(\text{C. }-\dfrac{1}{9}\)  & \(\text{D. }-\dfrac{7}{9}\)
        \end{tabular*}\par}
    \end{enumerate}
\end{description}

\begin{description}
    \item[二、选择题] 本题共4小题，每小题5分，共20分．在每小题给出的选项中，有多项符合题目要求．全部选对的得5分，部分选对的得2分，有选错的得0分．
    \begin{enumerate}[leftmargin=5pt]
      \addtocounter{enumi}{+8}
        \item 有一组样本数据 $x_1,x_2,\cdots,x_6$，其中 $x_1$ 是最小值，$x_6$ 是最大值，则
        \begin{enumerate}[label=\Alph*.]
            \item $x_2,x_3,x_4,x_5$ 的平均数等于 $x_1,x_2,\cdots,x_6$ 的平均数
            \item $x_2,x_3,x_4,x_5$ 的中位数等于 $x_1,x_2,\cdots,x_6$ 的中位数
            \item $x_2,x_3,x_4,x_5$ 的标准差不小于 $x_1,x_2,\cdots,x_6$ 的标准差
            \item $x_2,x_3,x_4,x_5$ 的极差不大于 $x_1,x_2,\cdots,x_6$ 的极差
        \end{enumerate}

        \item 噪声污染问题越来越受到重视．用声压级来度量声音的强弱，定义声压级 $L_p=20\times\lg\dfrac{p}{p_0}$，其中常数 $p_0\,(p_0>0)$ 是听觉下限阈值，$p$ 是实际声压．下表为不同声源的声压级：
        
        \begin{center}
        \begin{tabular}{c|c|c|c}
        \hline
        声源 & 与声源的距离/m & 声压级/dB \\
        \hline
        燃油汽车 & 10 & $60\sim90$ \\
        混合动力汽车 & 10 & $50\sim60$ \\
        电动汽车 & 10 & $40$ \\
        \hline
        \end{tabular}
        \end{center}

        已知在距离燃油汽车、混合动力汽车、电动汽车 $10\,\text{m}$ 处测得实际声压分别为 $p_1,p_2,p_3$，则
        \begin{enumerate}[label=\Alph*.]
            \item $p_1\geqslant p_2$
            \item $p_2>10p_3$
            \item $p_3=100p_0$
            \item $p_1\leqslant100p_2$
        \end{enumerate}

        \item 已知函数 $f(x)$ 的定义域为 $\mathbb{R}$，$f(xy)=y^2f(x)+x^2f(y)$，则
        \begin{enumerate}[label=\Alph*.]
            \item $f(0)=0$
            \item $f(1)=0$
            \item $f(x)$ 是偶函数
            \item $x=0$ 为 $f(x)$ 的极小值点
        \end{enumerate}

        \item 下列物体中，能够被整体放入棱长为 $1$（单位：m）的正方体容器（容器壁厚度忽略不计）内的有
        \begin{enumerate}[label=\Alph*.]
            \item 直径为 $0.99\,\text{m}$ 的球体
            \item 所有棱长均为 $1.4\,\text{m}$ 的四面体
            \item 底面直径为 $0.01\,\text{m}$，高为 $1.8\,\text{m}$ 的圆柱体
            \item 底面直径为 $1.2\,\text{m}$，高为 $0.01\,\text{m}$ 的圆柱体
        \end{enumerate}
    \end{enumerate}
\end{description}

\begin{description}
    \item[三、填空题] 本题共4小题，每小题5分，共20分．
    \begin{enumerate}[leftmargin=5pt]
      \addtocounter{enumi}{+12}
        \item 某学校开设了4门体育类选修课和4门艺术类选修课，学生需从这8门课中选修2门或3门课，并且每类选修课至少选修1门，则不同的选课方案共有 \rule[-0.3ex]{3em}{0.4pt} 种（用数字作答）．
        \item 在正四棱台 $ABCD$-$A_1B_1C_1D_1$ 中，$AB=2$，$A_1B_1=1$，$AA_1=\sqrt{2}$，则该棱台的体积为 \rule[-0.3ex]{3em}{0.4pt}．
        \item 已知函数 $f(x)=\cos\omega x-1\,(\omega>0)$ 在区间 $[0,2\pi]$ 有且仅有 $3$ 个零点，则 $\omega$ 的取值范围是 \rule[-0.3ex]{3em}{0.4pt}．
        \item 已知双曲线 $C:\dfrac{x^2}{a^2}-\dfrac{y^2}{b^2}=1\,(a>0,b>0)$ 的左、右焦点分别为 $F_1,F_2$．点 $A$ 在 $C$ 上，点 $B$ 在 $y$ 轴上，$\overrightarrow{F_1A}\perp\overrightarrow{F_1B}$，$\overrightarrow{F_2A}=-\dfrac{2}{3}\overrightarrow{F_2B}$，则 $C$ 的离心率为 \rule[-0.3ex]{3em}{0.4pt}．
    \end{enumerate}
\end{description}

\begin{description}
    \item[四、解答题] 本题共6小题，共70分．解答应写出文字说明、证明过程或演算步骤．
    \begin{enumerate}[leftmargin=5pt]
      \addtocounter{enumi}{+16}
        \item （10分）

        已知在 $\triangle ABC$ 中，$A+B=3C$，$2\sin(A-C)=\sin B$．

        \begin{enumerate}[label=(\arabic*)]
            \item 求 $\sin A$；
            \item 设 $AB=5$，求 $AB$ 边上的高．
        \end{enumerate}

        \item （12分）

        如图，在正四棱柱 $ABCD$-$A_1B_1C_1D_1$ 中，$AB=2$，$AA_1=4$．点 $A_2,B_2,C_2,D_2$ 分别在棱 $AA_1,BB_1,CC_1,DD_1$ 上，$AA_2=1$，$BB_2=DD_2=2$，$CC_2=3$．

        \textsf{[插入图片：正四棱柱示意图]}

        \begin{enumerate}[label=(\arabic*)]
            \item 证明：$B_2C_2\parallel A_2D_2$；
            \item 点 $P$ 在棱 $BB_1$ 上，当二面角 $P$-$A_2C_2$-$D_2$ 为 $150^\circ$ 时，求 $B_2P$．
        \end{enumerate}

        \item （12分）

        已知函数 $f(x)=a(\e^x+a)-x$．

        \begin{enumerate}[label=(\arabic*)]
            \item 讨论 $f(x)$ 的单调性；
            \item 证明：当 $a>0$ 时，$f(x)>2\ln a+\dfrac{3}{2}$．
        \end{enumerate}

        \item （12分）

        设等差数列 $\{a_n\}$ 的公差为 $d$，且 $d>1$．令 $b_n=\dfrac{n^2+n}{a_n}$，记 $S_n,T_n$ 分别为数列 $\{a_n\},\{b_n\}$ 的前 $n$ 项和．

        \begin{enumerate}[label=(\arabic*)]
            \item 若 $3a_2=3a_1+a_3$，$S_3+T_3=21$，求 $\{a_n\}$ 的通项公式；
            \item 若 $\{b_n\}$ 为等差数列，且 $S_{99}-T_{99}=99$，求 $d$．
        \end{enumerate}

        \item （12分）

        甲乙两人投篮，每次由其中一人投篮，规则如下：若命中则此人继续投篮，若末命中则换为对方投篮．无论之前投篮情况如何，甲每次投篮的命中率均为 $0.6$，乙每次投篮的命中率均为 $0.8$．由抽签确定第1次投篮的人选，第一次投篮的人是甲、乙的概率各为 $0.5$．

        \begin{enumerate}[label=(\arabic*)]
            \item 求第2次投篮的人是乙的概率；
            \item 求第 $i$ 次投篮的人是甲的概率；
            \item 已知：若随机变量 $X_i$ 服从两点分布，且 $P(X_i=1)=1-P(X_i=0)=q_i$，$i=1,2,\cdots,n$，则 $E\bigl(\sum_{i=1}^n X_i\bigr)=\sum_{i=1}^n q_i$．记前 $n$ 次（即从第1次到第 $n$ 次投篮）中甲投篮的次数为 $Y$，求 $E(Y)$．
        \end{enumerate}

        \item （12分）

        在直角坐标系 $xOy$ 中，点 $P$ 到 $x$ 轴的距离等于点 $P$ 到点 $\bigl(0,\frac{1}{2}\bigr)$ 的距离，记动点 $P$ 的轨迹为 $W$．

        \begin{enumerate}[label=(\arabic*)]
            \item 求 $W$ 的方程；
            \item 已知矩形 $ABCD$ 有三个顶点在 $W$ 上，证明：矩形 $ABCD$ 的周长大于 $3\sqrt{3}$．
        \end{enumerate}
    \end{enumerate}
\end{description}
\clearpage

\subsection{2023年新高考II卷}\label{gk:2023年新高考II卷}

\begin{description}
    \item[一、选择题] 本题共8小题，每小题5分，共40分．在每小题给出的四个选项中，只有一项是符合题目要求的．
    \begin{enumerate}[leftmargin=5pt]
        \item 在复平面内，$(1+3\i)(3-\i)$ 对应的点位于
       
        {\centering\begin{tabular*}{\linewidth}{@{}@{\extracolsep{\fill}}llll@{}}
          \(\text{A. }\)第一象限  & \(\text{B. }\)第二象限  & \(\text{C. }\)第三象限  & \(\text{D. }\)第四象限
        \end{tabular*}\par}

        \item 设集合 $A=\{0,-a\}$，$B=\{1,a-2,2a-2\}$，若 $A\subseteq B$，则 $a=$
       
        {\centering\begin{tabular*}{\linewidth}{@{}@{\extracolsep{\fill}}llll@{}}
          \(\text{A. }2\)  & \(\text{B. }1\)  & \(\text{C. }\dfrac{2}{3}\)  & \(\text{D. }-1\)
        \end{tabular*}\par}

        \item 某学校为了解学生参加体育运动的情况，用比例分配的分层随机抽样方法作抽样调查，拟从初中部和高中部两层共抽取60名学生，已知该校初中部和高中部分别有400名和200名学生，则不同的抽样结果共有
       
        {\centering\begin{tabular*}{\linewidth}{@{}@{\extracolsep{\fill}}llll@{}}
          \(\text{A. }\mathrm{C}_{400}^{45}\cdot\mathrm{C}_{200}^{15}\) 种  & \(\text{B. }\mathrm{C}_{400}^{20}\cdot\mathrm{C}_{200}^{40}\) 种 \\
          \(\text{C. }\mathrm{C}_{400}^{30}\cdot\mathrm{C}_{200}^{30}\) 种  & \(\text{D. }\mathrm{C}_{400}^{40}\cdot\mathrm{C}_{200}^{20}\) 种
        \end{tabular*}\par}

        \item 若 $f(x)=(x+a)\ln\dfrac{2x-1}{2x+1}$ 为偶函数，则 $a=$
       
        {\centering\begin{tabular*}{\linewidth}{@{}@{\extracolsep{\fill}}llll@{}}
          \(\text{A. }-1\)  & \(\text{B. }0\)  & \(\text{C. }\dfrac{1}{2}\)  & \(\text{D. }1\)
        \end{tabular*}\par}

        \item 已知椭圆 $C:\dfrac{x^2}{3}+y^2=1$ 的左、右焦点分别为 $F_1,F_2$，直线 $y=x+m$ 与 $C$ 交于 $A,B$ 两点，若 $\triangle F_1AB$ 面积是 $\triangle F_2AB$ 面积的 $2$ 倍，则 $m=$
       
        {\centering\begin{tabular*}{\linewidth}{@{}@{\extracolsep{\fill}}llll@{}}
          \(\text{A. }\dfrac{2}{3}\)  & \(\text{B. }\dfrac{\sqrt{2}}{3}\)  & \(\text{C. }-\dfrac{\sqrt{2}}{3}\)  & \(\text{D. }-\dfrac{2}{3}\)
        \end{tabular*}\par}

        \item 已知函数 $f(x)=a\e^x-\ln x$ 在区间 $(1,2)$ 单调递增，则 $a$ 的最小值为
       
        {\centering\begin{tabular*}{\linewidth}{@{}@{\extracolsep{\fill}}llll@{}}
          \(\text{A. }\e^2\)  & \(\text{B. }\e\)  & \(\text{C. }\e^{-1}\)  & \(\text{D. }\e^{-2}\)
        \end{tabular*}\par}

        \item 已知 $\alpha$ 为锐角，$\cos\alpha=\dfrac{1+\sqrt{5}}{4}$，则 $\sin\dfrac{\alpha}{2}=$
       
        {\centering\begin{tabular*}{\linewidth}{@{}@{\extracolsep{\fill}}llll@{}}
          \(\text{A. }\dfrac{3-\sqrt{5}}{8}\)  & \(\text{B. }\dfrac{-1+\sqrt{5}}{8}\)  & \(\text{C. }\dfrac{3-\sqrt{5}}{4}\)  & \(\text{D. }\dfrac{-1+\sqrt{5}}{4}\)
        \end{tabular*}\par}

        \item 记 $S_n$ 为等比数列 $\{a_n\}$ 的前 $n$ 项和，若 $S_4=-5$，$S_6=21S_2$，则 $S_8=$
       
        {\centering\begin{tabular*}{\linewidth}{@{}@{\extracolsep{\fill}}llll@{}}
          \(\text{A. }120\)  & \(\text{B. }85\)  & \(\text{C. }-85\)  & \(\text{D. }-120\)
        \end{tabular*}\par}
    \end{enumerate}
\end{description}

\begin{description}
    \item[二、选择题] 本题共4小题，每小题5分，共20分．在每小题给出的选项中，有多项符合题目要求．全部选对的得5分，部分选对的得2分，有选错的得0分．
    \begin{enumerate}[leftmargin=5pt]
      \addtocounter{enumi}{+8}
        \item 已知圆锥的顶点为 $P$，底面圆心为 $O$，$AB$ 为底面直径，$\angle APB=120^\circ$，$PA=2$，点 $C$ 在底面圆周上，且二面角 $P$-$AC$-$O$ 为 $45^\circ$，则
        \begin{enumerate}[label=\Alph*.]
            \item 该圆锥的体积为 $\pi$
            \item 该圆锥的侧面积为 $4\sqrt{3}\pi$
            \item $AC=2\sqrt{2}$
            \item $\triangle PAC$ 的面积为 $\sqrt{3}$
        \end{enumerate}

        \item 设 $O$ 为坐标原点，直线 $y=-\sqrt{3}(x-1)$ 过抛物线 $C:y^2=2px\,(p>0)$ 的焦点，且与 $C$ 交于 $M,N$ 两点，$l$ 为 $C$ 的准线，则
        \begin{enumerate}[label=\Alph*.]
            \item $p=2$
            \item $|MN|=\dfrac{8}{3}$
            \item 以 $MN$ 为直径的圆与 $l$ 相切
            \item $\triangle OMN$ 为等腰三角形
        \end{enumerate}

        \item 若函数 $f(x)=a\ln x+\dfrac{b}{x}+\dfrac{c}{x^2}\,(a\neq0)$ 既有极大值也有极小值，则
        \begin{enumerate}[label=\Alph*.]
            \item $bc>0$
            \item $ab>0$
            \item $b^2+8ac>0$
            \item $ac<0$
        \end{enumerate}

        \item 在信道内传输 $0,1$ 信号，信号的传输相互独立．发送 $0$ 时，收到 $1$ 的概率为 $\alpha\,(0<\alpha<1)$，收到 $0$ 的概率为 $1-\alpha$；发送 $1$ 时，收到 $0$ 的概率为 $\beta\,(0<\beta<1)$，收到 $1$ 的概率为 $1-\beta$．考虑两种传输方案：单次传输和三次传输．单次传输是指每个信号只发送 $1$ 次，三次传输是指每个信号重复发送 $3$ 次．收到的信号需要译码，译码规则如下：单次传输时，收到的信号即为译码；三次传输时，收到的信号中出现次数多的即为译码（例如，若依次收到 $1,0,1$，则译码为 $1$）．
        \begin{enumerate}[label=\Alph*.]
            \item 采用单次传输方案，若依次发送 $1,0,1$，则依次收到 $1,0,1$ 的概率为 $(1-\alpha)(1-\beta)^2$
            \item 采用三次传输方案，若发送 $1$，则依次收到 $1,0,1$ 的概率为 $\beta(1-\beta)^2$
            \item 采用三次传输方案，若发送 $1$，则译码为 $1$ 的概率为 $\beta(1-\beta)^2+(1-\beta)^3$
            \item 当 $0<\alpha<0.5$ 时，若发送 $0$，则采用三次传输方案译码为 $0$ 的概率大于采用单次传输方案译码为 $0$ 的概率
        \end{enumerate}
    \end{enumerate}
\end{description}

\begin{description}
    \item[三、填空题] 本题共4小题，每小题5分，共20分．
    \begin{enumerate}[leftmargin=5pt]
      \addtocounter{enumi}{+12}
        \item 已知向量 $\bm a,\bm b$ 满足 $|\bm a-\bm b|=\sqrt{3}$，$|\bm a+\bm b|=|2\bm a-\bm b|$，则 $|\bm b|=$ \rule[-0.3ex]{3em}{0.4pt}．
        \item 底面边长为 $4$ 的正四棱锥被平行于其底面的平面所截，截去一个底面边长为 $2$，高为 $3$ 的正四棱锥，所得棱台的体积为 \rule[-0.3ex]{3em}{0.4pt}．
        \item 已知直线 $x-my+1=0$ 与 $\odot C:(x-1)^2+y^2=4$ 交于 $A,B$ 两点，写出满足"$\triangle ABC$ 面积为 $\dfrac{8}{5}$"的 $m$ 的一个值 \rule[-0.3ex]{3em}{0.4pt}．
        \item 已知函数 $f(x)=\sin(\omega x+\varphi)$，如图，$A,B$ 是直线 $y=\dfrac{1}{2}$ 与曲线 $y=f(x)$ 的两个交点，若 $|AB|=\dfrac{\pi}{6}$，则 $f(\pi)=$ \rule[-0.3ex]{3em}{0.4pt}．
        
        \textsf{[插入图片：三角函数图象]}
    \end{enumerate}
\end{description}

\begin{description}
    \item[四、解答题] 本题共6小题，共70分．解答应写出文字说明、证明过程或演算步骤．
    \begin{enumerate}[leftmargin=5pt]
      \addtocounter{enumi}{+16}
        \item （10分）

        记 $\triangle ABC$ 的内角 $A,B,C$ 的对边分别为 $a,b,c$，已知 $\triangle ABC$ 面积为 $\sqrt{3}$，$D$ 为 $BC$ 的中点，且 $AD=1$．

        \begin{enumerate}[label=(\arabic*)]
            \item 若 $\angle ADC=\dfrac{\pi}{3}$，求 $\tan B$；
            \item 若 $b^2+c^2=8$，求 $b,c$．
        \end{enumerate}

        \item （12分）

        已知 $\{a_n\}$ 为等差数列，$b_n=\begin{cases}a_n-6,&n\text{为奇数},\\2a_n,&n\text{为偶数}.\end{cases}$ 记 $S_n,T_n$ 分别为数列 $\{a_n\},\{b_n\}$ 的前 $n$ 项和，$S_4=32$，$T_3=16$．

        \begin{enumerate}[label=(\arabic*)]
            \item 求 $\{a_n\}$ 的通项公式；
            \item 证明：当 $n>5$ 时，$T_n>S_n$．
        \end{enumerate}

        \item （12分）

        某研究小组经过研究发现某种疾病的患病者与未患病者的某项医学指标有明显差异，经过大量调查，得到如下的患病者和未患病者该指标的频率分布直方图：

        \textsf{[插入图片：患病者和未患病者频率分布直方图]}

        利用该指标制定一个检测标准，需要确定临界值 $c$，将该指标大于 $c$ 的人判定为阳性，小于或等于 $c$ 的人判定为阴性．此检测标准的漏诊率是将患病者判定为阴性的概率，记为 $p(c)$；误诊率是将未患病者判定为阳性的概率，记为 $q(c)$．假设数据在组内均匀分布．以事件发生的频率作为相应事件发生的概率．

        \begin{enumerate}[label=(\arabic*)]
            \item 当漏诊率 $p(c)=0.5\%$ 时，求临界值 $c$ 和误诊率 $q(c)$；
            \item 设函数 $f(c)=p(c)+q(c)$．当 $c\in[95,105]$，求 $f(c)$ 的解析式，并求 $f(c)$ 在区间 $[95,105]$ 的最小值．
        \end{enumerate}

        \item （12分）

        如图，三棱锥 $A$-$BCD$ 中，$DA=DB=DC$，$BD\perp CD$，$\angle ADB=\angle ADC=60^\circ$，$E$ 为 $BC$ 的中点．

        \textsf{[插入图片：三棱锥示意图]}

        \begin{enumerate}[label=(\arabic*)]
            \item 证明：$BC\perp DA$；
            \item 点 $F$ 满足 $\overrightarrow{EF}=\overrightarrow{DA}$，求二面角 $D$-$AB$-$F$ 的正弦值．
        \end{enumerate}

        \item （12分）

        已知双曲线 $C$ 的中心为坐标原点，左焦点为 $(-2\sqrt{5},0)$，离心率为 $\sqrt{5}$．

        \begin{enumerate}[label=(\arabic*)]
            \item 求 $C$ 的方程；
            \item 记 $C$ 的左、右顶点分别为 $A_1,A_2$，过点 $(-4,0)$ 的直线与 $C$ 的左支交于 $M,N$ 两点，$M$ 在第二象限，直线 $MA_1$ 与直线 $NA_2$ 交于 $P$，证明：点 $P$ 在定直线上．
        \end{enumerate}

        \item （12分）

        \begin{enumerate}[label=(\arabic*)]
            \item 证明：当 $0<x<1$ 时，$x-x^2<\sin x<x$；
            \item 已知函数 $f(x)=\cos ax-\ln(1-x^2)$，若 $x=0$ 是 $f(x)$ 的极大值点，求 $a$ 的取值范围．
        \end{enumerate}
    \end{enumerate}
\end{description}
\clearpage

\subsection{2023年北京卷}\label{gk:2023年北京卷}

\begin{description}
    \item[一、选择题] 本题共10小题，每小题4分，共40分．在每小题给出的四个选项中，选出符合题目要求的一项．
    \begin{enumerate}[leftmargin=5pt]
        \item 已知集合 $M=\{x\mid x+2\geqslant0\}$，$N=\{x\mid x-1<0\}$，则 $M\cap N=$
       
        {\centering\begin{tabular*}{\linewidth}{@{}@{\extracolsep{\fill}}llll@{}}
          \(\text{A. }\{x\mid -2\leqslant x<1\}\)  & \(\text{B. }\{x\mid -2<x\leqslant1\}\)  & \(\text{C. }\{x\mid x\geqslant-2\}\)  & \(\text{D. }\{x\mid x<1\}\)
        \end{tabular*}\par}

        \item 在复平面内，复数 $z$ 对应的点的坐标是 $(-1,\sqrt{3})$，则 $z$ 的共轭复数 $\bar{z}=$
       
        {\centering\begin{tabular*}{\linewidth}{@{}@{\extracolsep{\fill}}llll@{}}
          \(\text{A. }1+\sqrt{3}\i\)  & \(\text{B. }1-\sqrt{3}\i\)  & \(\text{C. }-1+\sqrt{3}\i\)  & \(\text{D. }-1-\sqrt{3}\i\)
        \end{tabular*}\par}

        \item 已知向量 $\bm a,\bm b$ 满足 $\bm a+\bm b=(2,3)$，$\bm a-\bm b=(-2,1)$，则 $|\bm a|^2-|\bm b|^2=$
       
        {\centering\begin{tabular*}{\linewidth}{@{}@{\extracolsep{\fill}}llll@{}}
          \(\text{A. }-2\)  & \(\text{B. }-1\)  & \(\text{C. }0\)  & \(\text{D. }1\)
        \end{tabular*}\par}

        \item 下列函数中，在区间 $(0,+\infty)$ 上单调递增的是
       
        {\centering\begin{tabular*}{\linewidth}{@{}@{\extracolsep{\fill}}llll@{}}
          \(\text{A. }f(x)=-\ln x\)  & \(\text{B. }f(x)=\dfrac{1}{2^x}\)  & \(\text{C. }f(x)=-\dfrac{1}{x}\)  & \(\text{D. }f(x)=3^{|x-1|}\)
        \end{tabular*}\par}

        \item 在 $\bigl(2x-\frac{1}{x}\bigr)^5$ 的展开式中，$x$ 的系数为
       
        {\centering\begin{tabular*}{\linewidth}{@{}@{\extracolsep{\fill}}llll@{}}
          \(\text{A. }-40\)  & \(\text{B. }40\)  & \(\text{C. }-80\)  & \(\text{D. }80\)
        \end{tabular*}\par}

        \item 已知抛物线 $C:y^2=8x$ 的焦点为 $F$，点 $M$ 在 $C$ 上．若 $M$ 到直线 $x=-3$ 的距离为 $5$，则 $|MF|=$
       
        {\centering\begin{tabular*}{\linewidth}{@{}@{\extracolsep{\fill}}llll@{}}
          \(\text{A. }7\)  & \(\text{B. }6\)  & \(\text{C. }5\)  & \(\text{D. }4\)
        \end{tabular*}\par}

        \item 在 $\triangle ABC$ 中，$(a+c)(\sin A-\sin C)=b(\sin A-\sin B)$，则 $\angle C=$
       
        {\centering\begin{tabular*}{\linewidth}{@{}@{\extracolsep{\fill}}llll@{}}
          \(\text{A. }\dfrac{\pi}{6}\)  & \(\text{B. }\dfrac{\pi}{3}\)  & \(\text{C. }\dfrac{2\pi}{3}\)  & \(\text{D. }\dfrac{5\pi}{6}\)
        \end{tabular*}\par}

        \item 若 $xy\neq0$，则"$x+y=0$"是"$\dfrac{y}{x}+\dfrac{x}{y}=-2$"的
       
        {\centering\begin{tabular*}{\linewidth}{@{}@{\extracolsep{\fill}}llll@{}}
          \(\text{A. }\)充分不必要条件  & \(\text{B. }\)必要不充分条件  & \(\text{C. }\)充要条件  & \(\text{D. }\)既不充分也不必要条件
        \end{tabular*}\par}

        \item 坡屋顶是我国传统建筑造型之一，蕴含着丰富的数学元素．安装灯带可以勾勒出建筑轮廓，展现造型之美．如图，某坡屋顶可视为一个五面体，其中两个面是全等的等腰梯形，两个面是全等的等腰三角形．若 $AB=25\,\text{m}$，$BC=10\,\text{m}$，且等腰梯形所在平面、等腰三角形所在平面与平面 $ABCD$ 的夹角的正切值均为 $\dfrac{\sqrt{14}}{5}$，则该五面体的所有棱长之和为
        
        \textsf{[插入图片：坡屋顶五面体示意图]}

        {\centering\begin{tabular*}{\linewidth}{@{}@{\extracolsep{\fill}}llll@{}}
          \(\text{A. }102\,\text{m}\)  & \(\text{B. }112\,\text{m}\)  & \(\text{C. }117\,\text{m}\)  & \(\text{D. }125\,\text{m}\)
        \end{tabular*}\par}

        \item 已知数列 $\{a_n\}$ 满足 $a_{n+1}=\dfrac{1}{4}(a_n-6)^3+6\,(n=1,2,3,\cdots)$，则
       
        {\centering\begin{tabular*}{\linewidth}{@{}@{\extracolsep{\fill}}llll@{}}
          \(\text{A. }\)当 $a_1=3$ 时，$\{a_n\}$ 为递减数列，且存在常数 $M\leqslant0$，使得 $a_n>M$ 恒成立 \\
          \(\text{B. }\)当 $a_1=5$ 时，$\{a_n\}$ 为递增数列，且存在常数 $M\leqslant6$，使得 $a_n<M$ 恒成立 \\
          \(\text{C. }\)当 $a_1=7$ 时，$\{a_n\}$ 为递减数列，且存在常数 $M>6$，使得 $a_n>M$ 恒成立 \\
          \(\text{D. }\)当 $a_1=9$ 时，$\{a_n\}$ 为递增数列，且存在常数 $M>0$，使得 $a_n<M$ 恒成立
        \end{tabular*}\par}
    \end{enumerate}
\end{description}

\begin{description}
    \item[二、填空题] 本题共5小题，每小题5分，共25分．
    \begin{enumerate}[leftmargin=5pt]
      \addtocounter{enumi}{+10}
        \item 已知函数 $f(x)=4^x+\log_2 x$，则 $f\bigl(\frac{1}{2}\bigr)=$ \rule[-0.3ex]{3em}{0.4pt}．
        \item 已知双曲线 $C$ 的焦点为 $(-2,0)$ 和 $(2,0)$，离心率为 $\sqrt{2}$，则 $C$ 的方程为 \rule[-0.3ex]{3em}{0.4pt}．
        \item 已知命题 $p$：若 $\alpha,\beta$ 为第一象限角，且 $\alpha>\beta$，则 $\tan\alpha>\tan\beta$．能说明 $p$ 为假命题的一组 $\alpha,\beta$ 的值为 $\alpha=$ \rule[-0.3ex]{3em}{0.4pt}，$\beta=$ \rule[-0.3ex]{3em}{0.4pt}．
        \item 我国度量衡的发展有着悠久的历史，战国时期就已经出现了类似于砝码的、用来测量物体质量的"环权"．已知9枚环权的质量（单位：铢）从小到大构成项数为9的数列 $\{a_n\}$，该数列的前3项成等差数列，后7项成等比数列，且 $a_1=1$，$a_5=12$，$a_9=192$，则 $a_7=$ \rule[-0.3ex]{3em}{0.4pt}；数列 $\{a_n\}$ 所有项的和为 \rule[-0.3ex]{3em}{0.4pt}．
        \item 设 $a>0$，函数 $f(x)=\begin{cases}x+2,&x<-a,\\\sqrt{a^2-x^2},&-a\leqslant x\leqslant a,\\-\sqrt{x}-1,&x>a.\end{cases}$ 给出下列四个结论：

        \ding{172}$f(x)$ 在区间 $(a-1,+\infty)$ 上单调递减；
        
         \ding{173}当 $a\geqslant1$ 时，$f(x)$ 存在最大值；
        
         \ding{174}设 $M(x_1,f(x_1))\,(x_1\leqslant a)$，$N(x_2,f(x_2))\,(x_2>a)$，则 $|MN|>1$；
        
         \ding{175}设 $P(x_3,f(x_3))\,(x_3<-a)$，$Q(x_4,f(x_4))\,(x_4\geqslant-a)$．若 $|PQ|$ 存在最小值，则 $a$ 的取值范围是 $\bigl(0,\frac{1}{2}\bigr]$．
        
        其中所有正确结论的序号是 \rule[-0.3ex]{3em}{0.4pt}．
    \end{enumerate}
\end{description}

\begin{description}
    \item[三、解答题] 本题共6小题，共85分．解答应写出文字说明、证明过程或演算步骤．
    \begin{enumerate}[leftmargin=5pt]
      \addtocounter{enumi}{+15}
        \item （13分）

        如图，在三棱锥 $P$-$ABC$ 中，$PA\perp$ 平面 $ABC$，$PA=AB=BC=1$，$PC=\sqrt{3}$．

        \textsf{[插入图片：三棱锥示意图]}

        \begin{enumerate}[label=(\arabic*)]
            \item 求证：$BC\perp$ 平面 $PAB$；
            \item 求二面角 $A$-$PC$-$B$ 的大小．
        \end{enumerate}

        \item （14分）

        设函数 $f(x)=\sin\omega x\cos\varphi+\cos\omega x\sin\varphi\,\bigl(\omega>0,|\varphi|<\frac{\pi}{2}\bigr)$．

        \begin{enumerate}[label=(\arabic*)]
            \item 若 $f(0)=-\dfrac{\sqrt{3}}{2}$，求 $\varphi$ 的值；
            \item 已知 $f(x)$ 在区间 $\bigl[-\frac{\pi}{3},\frac{2\pi}{3}\bigr]$ 上单调递增，$f\bigl(\frac{2\pi}{3}\bigr)=1$，再从条件 \ding{172}、条件 \ding{173}、条件 \ding{174}这三个条件中选择一个作为已知，使函数 $f(x)$ 的解析式唯一确定，求 $\omega$ 的值．
            \begin{quote}
            条件 \ding{172}：$x=\dfrac{\pi}{3}$ 是 $f(x)$ 的一条对称轴；
            
            条件 \ding{173}：$f\bigl(\dfrac{\pi}{12}\bigr)=\dfrac{\sqrt{3}}{2}$；
            
            条件 \ding{174}：$f\bigl(-\dfrac{\pi}{2}\bigr)=f\bigl(\dfrac{\pi}{2}\bigr)$．
            \end{quote}
        \end{enumerate}

        \item （14分）

        为研究某种农产品价格变化的规律，收集得到了该农产品连续40天的价格变化数据（如下表所示）．在描述价格变化时，用"$+$"表示"上涨"，即当天价格比前一天价格高；用"$-$"表示"下跌"，即当天价格比前一天价格低；用"$0$"表示"不变"，即当天价格与前一天价格相同．

        \textsf{[插入图片：价格变化数据表]}

        用频率估计概率．

        \begin{enumerate}[label=(\arabic*)]
            \item 试估计该农产品价格"上涨"的概率；
            \item 假设该农产品每天的价格变化是相互独立的．在未来的三个交易日中，用 $X$ 表示"上涨"的天数，求 $X$ 的分布列和数学期望；
            \item 用"上涨"天数占三个交易日的比例来估计该农产品价格的变化趋势，写出你的判断结论（不需要证明）．
        \end{enumerate}

        \item （15分）

        已知椭圆 $E:\dfrac{x^2}{a^2}+\dfrac{y^2}{b^2}=1\,(a>b>0)$ 的离心率为 $\dfrac{\sqrt{5}}{3}$，$A,C$ 分别是 $E$ 的上、下顶点，$B,D$ 分别是 $E$ 的左、右顶点，$|AC|=4$．

        \begin{enumerate}[label=(\arabic*)]
            \item 求 $E$ 的方程；
            \item 设 $P$ 为第一象限内 $E$ 上的动点，直线 $PD$ 与直线 $BC$ 交于点 $M$，直线 $PA$ 与直线 $y=-2$ 交于点 $N$．求证：$MN\parallel CD$．
        \end{enumerate}

        \item （15分）

        设函数 $f(x)=x-x^3\e^{ax+b}$，曲线 $y=f(x)$ 在点 $(1,f(1))$ 处的切线方程为 $y=-x+1$．

        \begin{enumerate}[label=(\arabic*)]
            \item 求 $a,b$ 的值；
            \item 设 $g(x)=f'(x)$，求 $g(x)$ 的单调区间；
            \item 求 $f(x)$ 的极值点个数．
        \end{enumerate}

        \item （15分）

        已知数列 $\{a_n\},\{b_n\}$ 的项数均为 $m\,(m>2)$，且 $a_n,b_n\in\{1,2,\cdots,m\}$，$\{a_n\},\{b_n\}$ 的前 $n$ 项和分别为 $A_n,B_n$，并规定 $A_0=B_0=0$．对于 $k\in\{0,1,2,\cdots,m\}$，定义

        \[
        r_k=\max\{i\mid B_i\leqslant A_k,i\in\{0,1,2,\cdots,m\}\}.
        \]

        \begin{enumerate}[label=(\arabic*)]
            \item 若 $a_n=1,b_n=2$，$m=4$，直接写出 $r_0,r_1,r_2,r_3,r_4$ 的值；
            \item 若 $\{a_n\}$ 为递增数列，求证：存在 $p,q\in\mathbb{N}^*$，$p>q$，使得 $A_p=B_q$；
            \item 若 $r_k\geqslant k$ 对任意 $k\in\{0,1,2,\cdots,m\}$ 均成立，且 $r_m=m$，求证：$\{a_n\}$ 与 $\{b_n\}$ 均为等差数列．
        \end{enumerate}
    \end{enumerate}
\end{description}
\clearpage

\subsection{2023年上海卷}\label{gk:2023年上海卷}

\begin{description}
    \item[一、填空题] 本大题共有12题，满分54分，第1\textasciitilde6题每题4分，第7\textasciitilde12题每题5分．
    \begin{enumerate}[leftmargin=5pt]
        \item 不等式 $|x-2|<1$ 的解集为 \rule[-0.3ex]{3em}{0.4pt}．
        \item 已知向量 $\bm a=(-2,3)$，$\bm b=(1,2)$，则 $\bm a\cdot\bm b=$ \rule[-0.3ex]{3em}{0.4pt}．
        \item 已知首项为 $3$，公比为 $2$ 的等比数列，设该等比数列的前 $n$ 项和为 $S_n$，则 $\log_2(S_n+3)=$ \rule[-0.3ex]{3em}{0.4pt}．
        \item 已知 $\tan\alpha=3$，则 $\tan2\alpha=$ \rule[-0.3ex]{3em}{0.4pt}．
        \item 已知函数 $f(x)=\begin{cases}2^x,&x>0,\\1,&x\leqslant0,\end{cases}$ 则 $f(x)$ 的值域为 \rule[-0.3ex]{3em}{0.4pt}．
        \item 已知当 $z=1+\i$，则 $|1-\i\cdot z|=$ \rule[-0.3ex]{3em}{0.4pt}．
        \item 已知圆 $x^2+y^2-4x-my=0$ 的面积为 $\pi$，则 $m=$ \rule[-0.3ex]{3em}{0.4pt}．
        \item 已知 $\triangle ABC$ 中，角 $A,B,C$ 所对的边 $a=4,b=5,c=6$，则 $\sin A=$ \rule[-0.3ex]{3em}{0.4pt}．
        \item 现有某地一年四个季度的 GDP（亿元），第一季度 GDP 为 $232$，第四季度 GDP 为 $241$，且四个季度的 GDP 的中位数为 $y$，平均数为 $\bar{y}$，若 $y-\bar{y}=1$，则第二、三季度的 GDP（从小到大排列）分别为 \rule[-0.3ex]{3em}{0.4pt}，\rule[-0.3ex]{3em}{0.4pt}．
        \item 已知 $(1+2023x)^{100}+(2023-x)^{100}=a_0+a_1x+a_2x^2+\cdots+a_{100}x^{100}$，若 $a_k<0$，$k\in\{0,1,2,\cdots,99\}$，则 $k=$ \rule[-0.3ex]{3em}{0.4pt}．
        \item 某公园修建一条新的健身步道，根据设计方案，步道有两个相互垂直的直线段构成，其长度为 $130\,\text{m}$，步道的起点和终点均为 $A$，将步道放置在平面直角坐标系 $xOy$ 中，$A$ 为原点，步道中不在坐标轴上的拐点为 $P$，在线段 $OP$ 中随机选取一点 $M$，在 $MP$ 中随机选取一点 $Q$，则称 $Q$ 为"步道内的点"．则"步道内的点"$Q$ 到点 $P$ 的距离小于 $40\,\text{m}$ 的概率为 \rule[-0.3ex]{3em}{0.4pt}．
        \item 已知 $a_1,a_2,a_3,a_4,a_5,a_6$ 是 $1,2,\cdots,6$ 的一个排列，满足 $a_1=2$，$a_2\neq1$，$a_3\neq3$，$a_4\neq4$，$a_5\neq5$，$a_6\neq6$，则满足条件的排列个数为 \rule[-0.3ex]{3em}{0.4pt}．
    \end{enumerate}
\end{description}

\begin{description}
    \item[二、选择题] 本题共4小题，每小题5分，共20分．
    \begin{enumerate}[leftmargin=5pt]
      \addtocounter{enumi}{+12}
        \item 要研究高一、高二、高三学生的平均身高，按性别采用分层抽样比较合适，以下说法正确的是
       
        {\centering\begin{tabular*}{\linewidth}{@{}@{\extracolsep{\fill}}llll@{}}
          \(\text{A. }\)简单随机抽样更合理  & \(\text{B. }\)按性别分层抽样合理  & \(\text{C. }\)按年级分层抽样合理  & \(\text{D. }\)两种分层抽样无差异
        \end{tabular*}\par}

        \item 下列函数中，在定义域内既是奇函数又是增函数的是
        
        {\centering\begin{tabular*}{\linewidth}{@{}@{\extracolsep{\fill}}llll@{}}
          \(\text{A. }y=x^3\)  & \(\text{B. }y=\tan x\)  & \(\text{C. }y=-\dfrac{1}{x}\)  & \(\text{D. }y=x|x|\)
        \end{tabular*}\par}

        \item 如图，在正方体 $ABCD$-$A_1B_1C_1D_1$ 中，$P$ 是 $A_1C_1$ 上的动点，则
        
        \textsf{[插入图片：正方体示意图]}

        {\centering\begin{tabular*}{\linewidth}{@{}@{\extracolsep{\fill}}llll@{}}
          \(\text{A. }\)存在点 $P$，使得 $AP\perp BC$ \\
          \(\text{B. }\)存在点 $P$，使得 $AP\perp BD_1$ \\
          \(\text{C. }\)对任意点 $P$，都有 $AP\parallel$ 平面 $A_1C_1D$ \\
          \(\text{D. }\)对任意点 $P$，都有 $AP\perp$ 平面 $BB_1D_1D$
        \end{tabular*}\par}

        \item 设 $f(x)=|x-a|+|x^3-ax^2|$，$x\in[0,1]$，$f(x)$ 的最大值为 $M(a)$，若 $M(a)$ 存在最小值，则 $a$ 的取值范围是
        
        {\centering\begin{tabular*}{\linewidth}{@{}@{\extracolsep{\fill}}llll@{}}
          \(\text{A. }(-\infty,0]\)  & \(\text{B. }[0,1]\)  & \(\text{C. }[1,+\infty)\)  & \(\text{D. }\mathbb{R}\)
        \end{tabular*}\par}
    \end{enumerate}
\end{description}

\begin{description}
    \item[三、解答题] 本题共5小题，共76分．
    \begin{enumerate}[leftmargin=5pt]
      \addtocounter{enumi}{+16}
        \item （14分）

        直四棱柱 $ABCD$-$A_1B_1C_1D_1$ 中，$AB\parallel DC$，$AB\perp AD$，$AB=2$，$AD=3$，$DC=4$．

        \textsf{[插入图片：直四棱柱示意图]}

        \begin{enumerate}[label=(\arabic*)]
            \item 求证：$A_1B\parallel$ 平面 $DCC_1D_1$；
            \item 若四棱柱 $ABCD$-$A_1B_1C_1D_1$ 的体积为 $36$，求二面角 $A_1$-$BD$-$A$ 的正切值．
        \end{enumerate}

        \item （14分）

        函数 $f(x)=\dfrac{x^2+3x+6}{x+1}\,(x>0)$．

        \begin{enumerate}[label=(\arabic*)]
            \item 求函数 $y=f(x)$ 的单调区间和极值；
            \item 若直线 $y=a$ 与 $y=f(x)$ 有 $3$ 个交点，求 $a$ 的取值范围．
        \end{enumerate}

        \item （14分）

        2023年6月7日，21世纪汽车博览会在上海举行．已知某汽车模型公司共有25个汽车模型，其外观和内饰的颜色分布如下表所示：

        \textsf{[插入表格：汽车模型颜色分布]}

        \begin{enumerate}[label=(\arabic*)]
            \item 若从25个模型中随机取一个，求取到红色外观模型的概率；
            \item 随机取两个模型，求两个模型外观颜色不同的概率；
            \item 该公司为答谢客户，对购买模型的客户进行抽奖活动，中奖的客户可以在三种奖品中任选一个．判断小红和小明在购买模型后选择相同奖品的概率是否大于 $\dfrac{1}{3}$，并说明理由．
        \end{enumerate}

        \item （17分）

        已知抛物线 $\Gamma:y^2=4x$，在 $\Gamma$ 上有一点 $A$ 位于第一象限，设 $A$ 的纵坐标为 $a\,(a>0)$．

        \begin{enumerate}[label=(\arabic*)]
            \item 若 $A$ 到 $\Gamma$ 的准线的距离为 $3$，求 $a$ 的值；
            \item 当 $a=4$ 时，若 $x$ 轴上存在一点 $B$，使 $AB$ 的中点在 $\Gamma$ 上，求 $O$ 到直线 $AB$ 的距离；
            \item 设直线 $l:x=-\sqrt{3}$，$P$ 是 $\Gamma$ 上异于 $A$ 的一个动点，$P$ 在 $l$ 上的投影为 $Q$，直线 $AP$ 与 $l$ 交于 $R$．若对任意点 $P$，$|QR|<4$，求 $a$ 的取值范围．
        \end{enumerate}

        \item （17分）

        已知 $f(x)=\ln x$，在 $y=f(x)$ 上取点 $P_1(a_1,\ln a_1)$，过 $P_1$ 作 $y=f(x)$ 的切线交 $x$ 轴于 $(a_2,0)$，过 $(a_2,0)$ 作 $x$ 轴的垂线交曲线 $y=f(x)$ 于 $P_2(a_2,\ln a_2)$，过 $P_2$ 作 $y=f(x)$ 的切线交 $x$ 轴于 $(a_3,0)$，如此继续．设 $a_1\in(1,+\infty)$，$a_n$ 构成数列 $\{a_n\}$．

        \begin{enumerate}[label=(\arabic*)]
            \item 若 $a_1=\e^2$，求 $a_2,a_3$；
            \item 求 $a_{n+1}$ 与 $a_n$ 的关系；
            \item 求证：$\displaystyle\lim_{n\to\infty}a_n=1$．
        \end{enumerate}
    \end{enumerate}
\end{description}
\clearpage

\subsection{2023年天津卷}\label{gk:2023年天津卷}

\begin{description}
    \item[一、选择题] 本题共9小题，每小题5分，共45分．在每小题给出的四个选项中，只有一项是符合题目要求的．
    \begin{enumerate}[leftmargin=5pt]
        \item 已知集合 $U=\{1,2,3,4,5\}$，$A=\{1,3\}$，$B=\{1,2,4\}$，则 $\complement_U B\cup A=$
       
        {\centering\begin{tabular*}{\linewidth}{@{}@{\extracolsep{\fill}}llll@{}}
          \(\text{A. }\{1,3,5\}\)  & \(\text{B. }\{1,3\}\)  & \(\text{C. }\{1,2,4\}\)  & \(\text{D. }\{1,2,4,5\}\)
        \end{tabular*}\par}

        \item "$a^2=b^2$"是"$a^2+b^2=2ab$"的
       
        {\centering\begin{tabular*}{\linewidth}{@{}@{\extracolsep{\fill}}llll@{}}
          \(\text{A. }\)充分不必要条件  & \(\text{B. }\)必要不充分条件  & \(\text{C. }\)充要条件  & \(\text{D. }\)既不充分也不必要条件
        \end{tabular*}\par}

        \item 若 $a=1.01^{0.5}$，$b=1.01^{0.6}$，$c=0.6^{0.5}$，则
       
        {\centering\begin{tabular*}{\linewidth}{@{}@{\extracolsep{\fill}}llll@{}}
          \(\text{A. }c>a>b\)  & \(\text{B. }c>b>a\)  & \(\text{C. }a>b>c\)  & \(\text{D. }b>a>c\)
        \end{tabular*}\par}

        \item 函数 $f(x)$ 的图象如下图所示，则 $f(x)$ 的解析式可能为
        
        \textsf{[插入图片：函数图象]}

        {\centering\begin{tabular*}{\linewidth}{@{}@{\extracolsep{\fill}}llll@{}}
          \(\text{A. }\dfrac{5(\e^x-\e^{-x})}{x^2+2}\)  & \(\text{B. }\dfrac{5\sin x}{x^2+1}\)  & \(\text{C. }\dfrac{5(\e^x+\e^{-x})}{x^2+2}\)  & \(\text{D. }\dfrac{5\cos x}{x^2+1}\)
        \end{tabular*}\par}

        \item 已知数列 $\{a_n\}$ 的前 $n$ 项和为 $S_n$，若 $a_1=2$，$a_{n+1}=2S_n+2\,(n\in\mathbb{N}^*)$，则 $a_4=$
       
        {\centering\begin{tabular*}{\linewidth}{@{}@{\extracolsep{\fill}}llll@{}}
          \(\text{A. }16\)  & \(\text{B. }32\)  & \(\text{C. }54\)  & \(\text{D. }162\)
        \end{tabular*}\par}

        \item 函数 $f(x)$ 的图象如下图所示，则 $f(x)$ 的解析式可能为
        
        \textsf{[插入图片：函数图象]}

        {\centering\begin{tabular*}{\linewidth}{@{}@{\extracolsep{\fill}}llll@{}}
          \(\text{A. }\dfrac{5(\e^x-\e^{-x})}{x^2+2}\)  & \(\text{B. }\dfrac{5\sin x}{x^2+1}\)  & \(\text{C. }\dfrac{5(\e^x+\e^{-x})}{x^2+2}\)  & \(\text{D. }\dfrac{5\cos x}{x^2+1}\)
        \end{tabular*}\par}

        \item 调查某种群花萼长度和花瓣长度，所得数据如图所示，其中相关系数 $r=0.8245$，下列说法正确的是
        
        \textsf{[插入图片：散点图]}

        {\centering\begin{tabular*}{\linewidth}{@{}@{\extracolsep{\fill}}llll@{}}
          \(\text{A. }\)花瓣长度和花萼长度没有相关性 \\
          \(\text{B. }\)花瓣长度和花萼长度呈现负相关 \\
          \(\text{C. }\)花瓣长度和花萼长度呈现正相关 \\
          \(\text{D. }\)若从样本中抽取一部分，则这部分的相关系数一定是 $0.8245$
        \end{tabular*}\par}

        \item 在三棱锥 $P$-$ABC$ 中，线段 $PC$ 上的点 $M$ 满足 $PM=\dfrac{1}{3}PC$，线段 $PB$ 上的点 $N$ 满足 $PN=\dfrac{2}{3}PB$，则三棱锥 $P$-$AMN$ 和三棱锥 $P$-$ABC$ 的体积之比为
       
        {\centering\begin{tabular*}{\linewidth}{@{}@{\extracolsep{\fill}}llll@{}}
          \(\text{A. }\dfrac{1}{9}\)  & \(\text{B. }\dfrac{2}{9}\)  & \(\text{C. }\dfrac{1}{3}\)  & \(\text{D. }\dfrac{4}{9}\)
        \end{tabular*}\par}

        \item 双曲线 $\dfrac{x^2}{a^2}-\dfrac{y^2}{b^2}=1\,(a>0,b>0)$ 的左、右焦点分别为 $F_1,F_2$．过 $F_2$ 作其中一条渐近线的垂线，垂足为 $P$．已知 $|PF_2|=2$，直线 $PF_1$ 的斜率为 $\dfrac{\sqrt{2}}{4}$，则双曲线的方程为
       
        {\centering\begin{tabular*}{\linewidth}{@{}@{\extracolsep{\fill}}llll@{}}
          \(\text{A. }\dfrac{x^2}{8}-\dfrac{y^2}{4}=1\)  & \(\text{B. }\dfrac{x^2}{4}-\dfrac{y^2}{8}=1\)  & \(\text{C. }\dfrac{x^2}{4}-\dfrac{y^2}{2}=1\)  & \(\text{D. }\dfrac{x^2}{2}-\dfrac{y^2}{4}=1\)
        \end{tabular*}\par}
    \end{enumerate}
\end{description}

\begin{description}
    \item[二、填空题] 本题共6小题，每小题5分，共30分．
    \begin{enumerate}[leftmargin=5pt]
      \addtocounter{enumi}{+9}
        \item 已知 $\i$ 是虚数单位，化简 $\dfrac{5+14\i}{2+3\i}$ 的结果为 \rule[-0.3ex]{3em}{0.4pt}．
        \item 在 $\bigl(2x^3-\dfrac{1}{x}\bigr)^6$ 的展开式中，常数项为 \rule[-0.3ex]{3em}{0.4pt}．
        \item 过原点的一条直线与圆 $C:(x+2)^2+y^2=3$ 相切，交曲线 $y^2=2px\,(p>0)$ 于点 $P$，若 $|OP|=8$，则 $p$ 的值为 \rule[-0.3ex]{3em}{0.4pt}．
        \item 甲乙丙三人投篮，每次由一人投篮，规则如下：第一次由甲投篮，若命中则此人继续投篮，若未命中则换为乙投篮；乙若命中则继续投篮，若未命中则换为丙投篮；丙若命中则继续投篮，若未命中则换为甲投篮，如此循环．设甲、乙、丙每次投篮的命中率分别为 $\dfrac{2}{3},\dfrac{1}{2},\dfrac{1}{3}$．则第2次投篮的人是乙的概率为 \rule[-0.3ex]{3em}{0.4pt}；记第 $n$ 次投篮的人是甲的概率为 $p_n$，则 $p_n=$ \rule[-0.3ex]{3em}{0.4pt}．
        \item 在 $\triangle ABC$ 中，$\angle A=60^\circ$，$BC=1$，点 $D$ 为 $AB$ 的中点，点 $E$ 为 $CD$ 的中点，若设 $\overrightarrow{AB}=\bm a$，$\overrightarrow{AC}=\bm b$，则 $\overrightarrow{AE}$ 可用 $\bm a,\bm b$ 表示为 \rule[-0.3ex]{3em}{0.4pt}；若 $\overrightarrow{BF}=\dfrac{1}{3}\overrightarrow{BC}$，则 $\overrightarrow{AE}\cdot\overrightarrow{AF}$ 的最大值为 \rule[-0.3ex]{3em}{0.4pt}．
        \item 若函数 $f(x)=ax^2-2x-|x^2-ax+1|$ 有且仅有两个零点，则 $a$ 的取值范围为 \rule[-0.3ex]{3em}{0.4pt}．
    \end{enumerate}
\end{description}

\begin{description}
    \item[三、解答题] 本题共5小题，共75分．解答应写出文字说明、证明过程或演算步骤．
    \begin{enumerate}[leftmargin=5pt]
      \addtocounter{enumi}{+15}
        \item （14分）

        在 $\triangle ABC$ 中，角 $A,B,C$ 所对的边分别为 $a,b,c$．已知 $a=\sqrt{39}$，$b=2$，$\angle A=120^\circ$．

        \begin{enumerate}[label=(\arabic*)]
            \item 求 $\sin B$ 的值；
            \item 求 $c$ 的值；
            \item 求 $\sin(B-C)$ 的值．
        \end{enumerate}

        \item （15分）

        如图，在三棱台 $ABC$-$A_1B_1C_1$ 中，已知 $A_1A\perp$ 平面 $ABC$，$AB\perp AC$，$AB=AC=AA_1=2$，$A_1C_1=1$，$N$ 为线段 $AB$ 的中点，$M$ 为线段 $BC$ 的中点．

        \textsf{[插入图片：三棱台示意图]}

        \begin{enumerate}[label=(\arabic*)]
            \item 求证：$A_1N\parallel$ 平面 $C_1MA$；
            \item 求平面 $C_1MA$ 与平面 $ACC_1A_1$ 所成角的余弦值；
            \item 求点 $C$ 到平面 $C_1MA$ 的距离．
        \end{enumerate}

        \item （15分）

        已知椭圆 $\dfrac{x^2}{a^2}+\dfrac{y^2}{b^2}=1\,(a>b>0)$ 的左、右顶点分别为 $A_1$ 和 $A_2$，右焦点为 $F$，且 $|A_1F|=3$，$|A_2F|=1$．

        \begin{enumerate}[label=(\arabic*)]
            \item 求椭圆的方程和离心率；
            \item 过 $A_1$ 的直线与椭圆交于点 $B$，垂直于 $x$ 轴的直线 $l$ 与 $A_2B$ 交于点 $T$，与 $x$ 轴交于点 $K$，且 $|TK|=|TA_2|$，求直线 $l$ 的方程．
        \end{enumerate}

        \item （15分）

        已知数列 $\{a_n\}$ 是等差数列，$a_2+a_5=16$，$a_5-a_3=4$．

        \begin{enumerate}[label=(\arabic*)]
            \item 求 $\{a_n\}$ 的通项公式和 $\displaystyle\sum_{i=2^{n-1}}^{2^n-1}a_i$；
            \item 已知 $\{b_n\}$ 为等比数列，对于任意 $k\in\mathbb{N}^*$，若 $2^{k-1}\leqslant n\leqslant2^k-1$，则 $b_k<a_n<b_{k+1}$．
            \begin{enumerate}[label=(\roman*)]
                \item 当 $k\geqslant2$ 时，求证：$2^k-1<b_k<2^k+1$；
                \item 求 $\{b_n\}$ 的通项公式及其前 $n$ 项和．
            \end{enumerate}
        \end{enumerate}

        \item （16分）

        已知函数 $f(x)=\bigl(\dfrac{1}{x}+\dfrac{1}{2}\bigr)\ln(x+1)$．

        \begin{enumerate}[label=(\arabic*)]
            \item 求曲线 $y=f(x)$ 的斜率为 $1$ 的切线的方程；
            \item 当 $x\in(0,+\infty)$ 时，求证：$f(x)>1$；
            \item 设 $F(x)=|f(x)-a|$，记 $F(x)$ 在区间 $[-1,1]$ 上的最大值为 $M(a)$．当 $M(a)$ 最小时，求 $a$ 的值．
        \end{enumerate}
    \end{enumerate}
\end{description}
\clearpage

